% Stable NG-01--NG-33 records, reorganized by mathematical cause.
% Historical review grades are preserved in the provenance record.
\section{How to read the negative results}
\label{S:nogo-guide}

The proof of the main results did not follow one route from beginning to end.
Several plausible steps lost the required time scale, changed the covariance
class, discarded part of the wave state, or asserted uniformity that was not
available. This chapter explains those failures as mathematical results.
Each entry starts with the reason the approach was attractive, states a
precise obstruction, and supplies its argument. It then separates what fails
from what survives and specifies what a replacement would need to establish.

\paragraph{Where to begin.}
For pairing arguments, start with the two-chain expansion in
Proposition~\ref{P:mixed}; for covariance and phase arguments, recall
Hypothesis~\ref{H:homog} and the displacement--velocity state. The topology
section states its own spaces, and the finite examples need only Gaussian
moments and elementary optimization. Each entry gives more specific
prerequisites and points to the part of the positive proof it explains.
The five groups below can be read independently.

The stable records NG-01 to NG-33 are printed below as No-go~01 to No-go~33.
Their numbers do not change when related examples are placed together. Some
restate results proved earlier in the book; others give the retained
counterexample in full. They are records of attempted approaches, not a count
of independent discoveries. No-go~\ref{P:nogo-33} groups three elementary
errors found in review arguments.

\paragraph{The scope of a counterexample.}
Unless an entry specifies a different model, the covariance and initial data
are those of Hypothesis~\ref{H:homog} and the canonical Stratonovich equation
with data $(1,0)$. We write $a=\alpha_0$. Scalar ordinary differential
equations, finite matrices and abstract Hilbert-space examples are identified
as such. They test the general principles stated in those entries; an example
outside Hypothesis~\ref{H:homog} does not refute a theorem restricted to that
hypothesis. A numerical majorant growing too quickly is an obstruction to its
use in a proof, not a lower bound on the actual moment.

\paragraph{The status of the evidence.}
The historical record distinguishes arguments already in the book, arguments
that received a review, and author-side diagnostics. These descriptions are
not interchangeable. A review may have occurred without meeting its required
acceptance threshold, as happened for the earlier NG-09 record. The proofs
below are the present mathematical account; the original reports and their
grades remain unchanged. Where an earlier computation was only formal, the
necessary analytic justification is given explicitly rather than attributed
to that earlier computation.

\begin{center}
\begin{tabular}{@{}>{\raggedright\arraybackslash}p{0.26\textwidth}>{\raggedright\arraybackslash}p{0.41\textwidth}>{\raggedright\arraybackslash}p{0.24\textwidth}@{}}
  \toprule
  Obstruction & What the book keeps explicit & Entries \\
  \midrule
  \hyperref[S:nogo-pairings]{Pairing and coefficient estimates} & Temporal order, spatial support and
  the quantifiers on constants & 01--07, 10, 17, 18 \\
  \hyperref[S:nogo-covariance]{Covariance comparisons} & The actual covariance class and the correct
  order of comparison & 08, 09, 11--13, 19 \\
  \hyperref[S:nogo-phase]{Phase and Gaussian controls} & Velocity, the second-order mean and the
  full tilted law & 14--16, 20, 31 \\
  \hyperref[S:nogo-topology]{Operator spaces and limits} & The specified topology, radius dependence
  and control of variational suprema & 21, 22, 24--27, 30, 32 \\
  \hyperref[S:nogo-optimization]{Finite optimization and review} & Both variances, active sets and exact
  algebraic factors & 23, 28, 29, 33 \\
  \bottomrule
\end{tabular}
\end{center}

The main text cites these entries where they explain a choice in the proof.
The entries also point back to the construction or estimate that survives.
This is intended to let a reader begin with an attempted approach, understand
its obstruction, and then follow the successful replacement.

Tao has observed that negative results are underreported, which can make the
public account of AI-assisted mathematics appear more successful than the
underlying research process~\cite{tao:26:post}*{first paragraph}. The purpose here is to make
these particular failures usable. Separate source-paper qualifications and
failures of the checking process retain their own records; they are not
silently added to the 33 mathematical identifiers below.

\begin{figure}[htbp]
  \centering
  % Vector metaphor only; the cards are the book's stable No-go identifiers.
\begin{tikzpicture}[x=1cm,y=1cm,font=\normalsize]
  \fill[hamblue!5] (-6.8,-4) rectangle (6.8,0);
  \filldraw[draw=hamblue!65,fill=hamblue!15,thick]
  (-1,0)--(-5.7,-.7)--(-6.1,-1.5)--(-6.2,-2.7)
  --(-4.8,-3.6)--(-1.2,-4)--(3.8,-3.75)
  --(5.8,-2.9)--(6.15,-1.5)--(5.7,-.7)--(1,0)--cycle;
  \filldraw[draw=hamblue,fill=white,thick]
  (-1,0)--(-.6,.85)--(-.05,1.85)--(.5,1.05)--(1,0)--cycle;
  \draw[hamblue,thick] (-4.4,0)--(6.8,0);
  \node[anchor=west,fill=white,inner sep=3pt] at (-6.6,.18) {waterline};
  \node[anchor=west,text=hamteal] at (2.15,1.35)
  {\ProofRef{Theorem}{T:intro-upper}};
  \node[anchor=west,text=hamteal] at (2.15,.65)
  {\ProofRef{Theorem}{T:intro-second-moment}};
  \draw[hamteal] (.55,1.05)--(1.5,1.35)--(2,1.35);
  \draw[hamteal] (.8,.45)--(1.5,.65)--(2,.65);
  \foreach \id/\xx in {01/-4.2,02/-1.4,33/4.2} {
  \node[draw=hamblue!70,fill=white,rounded corners=1.5pt,
  minimum width=1.9cm,minimum height=.55cm,inner sep=4pt]
  at (\xx,-1.95) {\hyperref[P:nogo-\id]{No-go~\id}};
  }
  \node at (1.4,-1.95) {$\cdots$};
  \node[font=\normalsize\bfseries,text=hamblue] at (0,-3.3)
  {Chapter 12: below the waterline};
\end{tikzpicture}

  \caption[What lies below the waterline]{The results stated in Chapter~1,
    Theorems~\ref{T:intro-upper} and~\ref{T:intro-second-moment}, are what shows
    above the waterline. This chapter is what lies below it: the 33 numbered
    obstructions and failed approaches, each kept with the reason it was tried and
    the argument that stops it.}
  \label{Fig:nogo-iceberg}
\end{figure}


\section{Pairings, ordering, and comparison kernels}\label{S:nogo-pairings}

The second moment contains two ordered chains and all their Gaussian
pairings, including pairs within a chain. This structure offers several
natural reductions: estimate each pairing separately, compare with a single
chain, or copy the first-moment clock. The obstructions below distinguish
loss of scale from failure of an inequality. A majorant can be too large even
when every step used to obtain it is valid.

Throughout this section write $a=\alpha_0$ and
$\beta=4-\alpha-a$. Let $X_n$ denote the ordered Stratonovich term of degree
$n$, and put
\[
  a_k=\E X_{2k}(1,x),\qquad
  A_k=\sum_{m+n=2k}\E[X_m(1,x)X_n(1,x)],\qquad
  b_1=\E X_1(1,x)^2.
\]
Thus $A_1=2a_1+b_1$. Homogeneity supplies the factor $t^{\beta k}$.
The nonnegative convergent moment expansions used here are those of
Proposition~\ref{P:mixed}; they include unequal chain lengths. In dimension
one we use $G_s(z)=\frac12\mathbf1_{\{|z|<s\}}$.

\subsection{NG-01: one common estimate loses a factorial power}
\label{S:nogo-01}

Theorem~3.1 of~\cite{chen:author2026:time-dependency} controls all the
pairing patterns. Applying its homogeneous norm bound to each pattern is
therefore an attractive first attempt at the upper exponent. The issue is
what happens after optimization and summation.

\setcounter{nogo}{0}
\begin{nogo}[The common-parameter majorant has the wrong scale]
  \label{P:nogo-01}
  Under Hypothesis~\ref{H:homog}, let $P_k=(2k)!/(2^k k!)$ and $K>0$.
  The numerical majorant
  \[
    B_K(t)=1+\sum_{k\ge1}(2k+1)P_kK^k
    \inf_{\theta>0}e^{2\theta t}\theta^{-\beta k}
  \]
  satisfies $\log B_K(t)\asymp t^{\beta/(\beta-1)}$ as $t\to\infty$.
  Since $\beta/(\beta-1)>\kappa=\beta/(3-\alpha)$, this majorant cannot
  establish a finite upper constant on the scale $t^\kappa$.
\end{nogo}
\begin{proof}
  The infimum equals $e^{\beta k}(2t/(\beta k))^{\beta k}$.
  The inequalities $k!\le P_k\le2^k k!$ and
  $e^{-k}\le k!/k^k\le1$ show that $B_K$ lies between two series
  \[
    \sum_{k\ge0}\frac{(c t^\beta)^k}{(k!)^q},\qquad q=\beta-1>1,
  \]
  with different positive constants $c$. For $z>0$, this series is at most
  $\exp(qz^{1/q})$, by the $\ell^q$--$\ell^1$ inequality. Its term with
  $k=\lfloor z^{1/q}/2\rfloor$ is at least $2^{qk}$ for large $z$.
  This proves the two-sided growth. Finally, $\beta-1=3-\alpha-a<3-\alpha$.
  This proves No-go~\ref{P:nogo-01}.
\end{proof}

This restates the obstruction in Proposition~\ref{P:envelope} and
Remark~\ref{R:G0-verdict}. The lower estimate is for $B_K$, not for the
moment. The valid upper estimate improves the crude quadratic exponent,
but its factorial denominator misses $(k!)^a$. The split count $2k+1$
costs only an exponential factor in $k$; removing it does not repair the
loss. The same calculation for one chain, with $e^{\theta t}$, has the
same defect. Reopening requires an estimate retaining ordering or grouping
pairings before taking this common bound. Section~\ref{S:interval-upper}
provides such a different argument.

\subsection{NG-02: a single adjacent diagram is too large}
\label{S:nogo-02}

There are roughly an exponential factor times $k!$ pairings of $2k$
vertices. In white space, a uniform denominator $(k!)^3$ for each diagram
would therefore give the desired denominator $(k!)^2$ after summation.
The following example explains why that convenient sufficient estimate is
stronger than the problem permits.

\setcounter{nogo}{1}
\begin{nogo}[No uniform cubic factorial denominator per diagram]
  \label{P:nogo-02}
  Let $d=1$, $\Gamma=\delta_0$ and $0<a<1$. For the split $(0,2k)$
  and adjacent matching $\{(1,2),(3,4),\ldots,(2k-1,2k)\}$, its contribution is
  \[
    J_{\rm adj}(t)=
    \left[\frac{\Gamma_{\mathrm E}(1-a)}2\right]^k
    \frac{t^{(3-a)k}}{\Gamma_{\mathrm E}(1+(3-a)k)}.
  \]
  Consequently, there is no finite $C=C(a)$ bounding every diagram by
  $C^k t^{(3-a)k}/(k!)^3$ for all $k$ and $t>0$.
\end{nogo}
\begin{proof}
  Use root-forward increments $h_1,\ldots,h_{2k}>0$ with
  $\sum h_i<t$. Each adjacent contraction evaluates the even Green factor
  at zero, giving $1/2$, and integration of the preceding free increment
  gives mass $h_{2j-1}$. The temporal factor is $h_{2j}^{-a}$. Therefore,
  \[
    J_{\rm adj}(t)=2^{-k}
    \int_{\sum h_i<t}\prod_{j=1}^k h_{2j-1}h_{2j}^{-a}\,\ud h.
  \]
  Here delta contractions mean limits of Gaussian approximate identities.
  For completeness, fix $k$ and a strictly positive increment vector.
  After convolution the paired positions approach equality. The even Green
  factors are continuous at zero; the remaining indicator boundaries are
  finitely many hyperplanes in the $k$ free positions and have measure zero.
  The limit is thus the stated diagonal integral. All positions in the
  original chain lie in $[-t,t]$. Uniformly in the smoothing widths, replacing
  its Green indicators by one bounds the spatial integral by
  $2^{-2k}(2t)^k$: each disjoint smoothed pair has integral at most $2t$
  over $[-t,t]^2$. The temporal factors on the full time cube are a product
  of $k$ disjoint factors $|s-r|^{-a}$, each integrable because $a<1$.
  Dominated convergence justifies the fixed-order spatial limit through the
  time integration, and also temporal truncation limits.

  The Dirichlet simplex formula, with parameters $2$ and $1-a$ for each
  pair, now gives the asserted expression. Stirling's formula yields
  \[
    \log\big((k!)^3J_{\rm adj}(1)\big)=a k\log k+O(k).
  \]
  This exceeds $k\log C$ eventually for every fixed $C$.
  This proves No-go~\ref{P:nogo-02}.
\end{proof}

The adjacent matching is only one diagram. Its denominator is, up to
exponential factors, $(k!)^{3-a}$, which is stronger than the desired
summed denominator $(k!)^2$. Hence, this witness does not obstruct the
summed bound. A viable successor must retain the number of large diagrams,
or sum before estimating; it cannot assign the same cubic denominator to
all of them.

\subsection[NG-03: a squared-mean comparison fails]{NG-03: doubled variance does not give the proposed squared mean}
\label{S:nogo-03}

A positive mixture of exponentials of linear Gaussian functionals would
satisfy a useful comparison between the second moment and the mean at
noise amplitude $\sqrt2$. Such a representation must be established for
the wave rather than inferred from a parabolic analogy.

\setcounter{nogo}{2}
\begin{nogo}[Failure of the all-time comparison with unit prefactor]
  \label{P:nogo-03}
  For the admissible Riesz case $d=1$, $\Gamma(x)=|x|^{-1/2}$ and $a=1/2$,
  \[
    \E u(t,x)^2>\big(\E u_{\sqrt2}(t,x)\big)^2
  \]
  for all sufficiently small $t>0$, where $u_\lambda$ uses noise amplitude
  $\lambda$ and the same initial data.
\end{nogo}
\begin{proof}
  More generally, for $\Gamma(x)=c|x|^{-\alpha}$, $0<\alpha<1$,
  direct integration of the length-two chain gives
  \[
    a_1=\frac{c}{(1-\alpha)(2-\alpha-a)(3-\alpha-a)(4-\alpha-a)}.
  \]
  Indeed the spatial integral is
  $c h_1h_2^{1-\alpha}/(1-\alpha)$ and its temporal factor is $h_2^{-a}$;
  integrate over $h_1+h_2<1$. At $\alpha=a=1/2$, $c=1$, this gives $a_1=1/3$.
  For independent uniforms $U$ on $[-s,s]$ and $V$ on $[-r,r]$, $0<r\le s$,
  \[
    \E|U-V|=\frac{s}{2}+\frac{r^2}{6s}\le\frac{2s}{3}.
  \]
  Jensen's inequality for $z\mapsto z^{-\alpha}$ implies
  \[
    \int G_s(x)G_r(y)|x-y|^{-\alpha}\,\ud x\ud y
    \ge (3/2)^\alpha s^{1-\alpha}r.
  \]
  Integrating on the two temporal triangles gives
  \[
    b_1\ge
    \frac{2(3/2)^\alpha}{(1-a)(2-a)(4-\alpha-a)}
    =\frac89\sqrt{\frac32}>\frac23=2a_1.
  \]
  The comparison is a finite-time failure, not just a formal coefficient test.
  Put $z=t^\beta$. The two nonnegative series under comparison have expansions
  \[
    \E u(t,x)^2=1+(2a_1+b_1)z+O(z^2),\qquad
    (\E u_{\sqrt2}(t,x))^2=1+4a_1z+O(z^2).
  \]
  To justify the remainders, choose any fixed positive $z_0$ at which these
  series converge, using the moment expansion also at amplitude $\sqrt2$.
  For each nonnegative coefficient sequence $(c_j)$,
  $\sum_{j\ge2}c_jz^j\le(z/z_0)^2\sum_{j\ge2}c_jz_0^j$ when
  $0\le z\le z_0$. Apply this to both series, including the nonnegative
  Cauchy product for the squared mean. Their difference therefore has the
  strictly positive leading coefficient $b_1-2a_1$.
  This proves No-go~\ref{P:nogo-03}.
\end{proof}

For a deterministic nonnegative mixture of Gaussian exponentials, the
inequality follows by Tonelli and
$\|h+k\|^2/2\le\|h\|^2+\|k\|^2$, provided the same mixture represents
all amplitudes in the comparison. That proposed mechanism is excluded by
the example. White noise, an eventual large-time comparison, or a prefactor
larger than one is not excluded. Reopening requires a restriction or a new
comparison that removes the contradicted first-coefficient requirement;
it cannot merely assert this positive-mixture representation.

\subsection[NG-04: constants near the time boundary]{NG-04: a coefficient constant cannot be uniform at the time boundary}
\label{S:nogo-04}

Comparing $A_k$ with $a_k$ retains the self-pairings of the Stratonovich
problem and would transfer a first-moment bound to the second moment.
However, a prescribed asymptotic constant must not be imposed on every
low-order coefficient without checking parameter uniformity.

\setcounter{nogo}{3}
\begin{nogo}[Divergence of the first coefficient ratio]
  \label{P:nogo-04}
  Fix $0<\alpha<1$ and $c>0$, with $\Gamma(x)=c|x|^{-\alpha}$ in $d=1$.
  Then
  \[
    \lim_{a\uparrow1}(1-a)\frac{A_1}{a_1}
    =2^{2-\alpha}(1-\alpha)>0.
  \]
  In particular no constant depending on $\alpha$ alone, including
  $2^{4-\alpha}$, can satisfy $A_k\le C^k a_k$ for every $k\ge1$
  and every $0<a<1$ at this fixed spatial covariance.
\end{nogo}
\begin{proof}
  Let
  $K(s,r)=\int G_s(x)G_r(y)c|x-y|^{-\alpha}\,\ud x\ud y$.
  This is continuous on $[0,1]^2$ by local integrability of the Riesz kernel
  and dominated convergence on a fixed square. Direct integration gives
  \[
    K(s,s)=\frac{c\,2^{1-\alpha}s^{2-\alpha}}
    {(1-\alpha)(2-\alpha)}.
  \]
  Writing $v=s-r$, define the continuous function
  $H(v)=2\int_v^1 K(s,s-v)\,\ud s$. Then
  \[
    b_1=\int_0^1v^{-a}H(v)\,\ud v,\qquad
    H(0)=\frac{c\,2^{2-\alpha}}
    {(1-\alpha)(2-\alpha)(3-\alpha)}.
  \]
  For any continuous $H$ on $[0,1]$,
  $\varepsilon\int_0^1v^{-1+\varepsilon}H(v)\,\ud v\to H(0)$.
  Indeed the measures have mass one, and their mass on $[\delta,1]$ is
  $1-\delta^\varepsilon\to0$; continuity controls $[0,\delta]$.
  Thus $(1-a)b_1\to H(0)$. The formula for $a_1$ in the preceding proof gives
  \[
    a_1\longrightarrow
    \frac{c}{(1-\alpha)^2(2-\alpha)(3-\alpha)}.
  \]
  Combine these limits with $A_1=2a_1+b_1$.
  This proves No-go~\ref{P:nogo-04}.
\end{proof}

For each fixed $a<1$, a finite constant $C=C(\alpha,a)$ remains possible;
the argument only forces it to diverge as $a\uparrow1$. It also says
nothing against comparisons beginning at large order or large time with
suitable prefactors. Reopening requires explicitly allowing that parameter
dependence or removing the low orders. Historical static-endpoint tests
likewise constrain an endpoint or a uniform family, not a separately
chosen constant at an arbitrary positive $a$.

\subsection{NG-05: chronological merging breaks wave support}
\label{S:nogo-05}

A single chain has the monotone ordering needed by the first-moment
representation. Sorting two chains by time looks like a direct way to
recover it. Finite propagation prevents the resulting pointwise comparison.

\setcounter{nogo}{4}
\begin{nogo}[Support obstruction to pointwise concatenation]
  \label{P:nogo-05}
  In $d=1$, no finite constant can bound every product of two root-forward
  chain kernels by the single-chain kernel obtained by chronological merging.
  The failure occurs in the $(2,2)$ split on a set of positive measure for
  Riesz covariance and, with domestic pairs, for white covariance.
\end{nogo}
\begin{proof}
  Work at horizon one with root $(0,0)$. The first chain has times
  $(1/4,3/4)$ and positions $(x,x)$, $x=1/5$; the second has times
  $(1/2,7/8)$ and positions $(y,y)$, $y=-2/5$.
  The four original displacements and elapsed times are
  \[
    (1/5,1/4),\quad(0,1/2),\quad(2/5,1/2),\quad(0,3/8).
  \]
  Each displacement is strictly smaller than its elapsed time, so their
  Green product is $2^{-4}$. The merged chain instead includes the jump
  from $(1/4,1/5)$ to $(1/2,-2/5)$. Its displacement is $3/5>1/4$,
  so that Green factor vanishes. All these inequalities are strict.

  For white noise choose domestic pairs within each original chain. The
  delta contractions leave precisely the two free spatial variables $x,y$.
  Open neighborhoods of the displayed four times and these two variables
  preserve both conclusions and have positive contracted measure. Gaussian
  regularization justifies this contraction as in No-go~\ref{P:nogo-02}.
  For Riesz covariance perturb all four spatial positions independently
  in small neighborhoods. Avoiding the domestic diagonals gives a set of
  positive ordinary measure on which the covariance factors are finite and
  strictly positive, while the same support defect persists.
  No multiplier can turn the zero merged product into a positive upper bound.
  This proves No-go~\ref{P:nogo-05}.
\end{proof}

\begin{figure}[htbp]
  \centering
  \begin{tikzpicture}[x=5cm,y=5cm,>=stealth,font=\small]
    \draw[->,gray] (-.65,0) -- (.55,0) node[right] {$x$};
    \draw[->,gray] (0,-.04) -- (0,1.04) node[above] {$s$};
    \draw[gray!35,densely dotted] (-.6,.6) -- (0,0) -- (.5,.5);
    \coordinate (O) at (0,0);
    \coordinate (A) at (.2,.25);
    \coordinate (B) at (.2,.75);
    \coordinate (C) at (-.4,.5);
    \coordinate (D) at (-.4,.875);
    \draw[hamblue,very thick,->] (O) -- (A);
    \draw[hamblue,very thick,->] (A) -- (B);
    \draw[hamteal,very thick,->] (O) -- (C);
    \draw[hamteal,very thick,->] (C) -- (D);
    \draw[hamgold,very thick,dashed,->] (A) -- (C)
    node[midway,above,sloped,fill=white,inner sep=2pt]
    {$3/5>1/4$};
    \fill (O) circle (.6pt) node[below right] {root $(0,0)$};
    \fill[hamblue] (A) circle (.7pt)
    node[right,black] {$(1/5,1/4)$};
    \fill[hamblue] (B) circle (.7pt)
    node[right,black] {$(1/5,3/4)$};
    \fill[hamteal] (C) circle (.7pt)
    node[left,black] {$(-2/5,1/2)$};
    \fill[hamteal] (D) circle (.7pt)
    node[left,black] {$(-2/5,7/8)$};
    \node[hamblue,anchor=west] at (.65,.94) {first chain};
    \node[hamteal,anchor=west] at (.65,.84) {second chain};
    \node[hamgold,anchor=west] at (.65,.74) {merged jump};
  \end{tikzpicture}
  \caption{Root-forward coordinates $(x,s)$ for No-go~\ref{P:nogo-05}.
    Each solid edge lies strictly inside its propagation cone. Chronological
    merging inserts the dashed jump, whose distance $3/5$ exceeds its elapsed
    time $1/4$. The picture refutes pointwise concatenation, not an integrated
    comparison. For domestic white pairs the two free positions are \mbox{$x$ and $y$}.}
  \label{F:nogo-chain-geometry}
\end{figure}


The time coordinates here are root-forward chain coordinates; applying a
terminal-to-point convention again would change the object being tested.
The result concerns pointwise concatenation. Integrated comparisons and
comparisons of grouped Wick sums remain possible. Reopening requires
averaging, extra displacement variables, or another mechanism that preserves
the two original propagation constraints instead of imposing the false
merged one.

\subsection{NG-06: erasing geometry before summation loses the scale}
\label{S:nogo-06}

After white contractions, the temporal pairing sum is symmetric. One can
exploit that symmetry immediately by forgetting the spatial indicators.
The resulting bound is easy to sum but too weak for the desired exponent.

\setcounter{nogo}{5}
\begin{nogo}[The geometry-free majorant]
  \label{P:nogo-06}
  For $d=1$, $\Gamma=\delta_0$, $0<a<1$, replacing every wave indicator by
  one while retaining only the spatial box $[-t,t]$ gives the coefficient
  majorant
  \[
    A_k t^{(3-a)k}\le
    \frac{1}{k!}\left[\frac{2t^{3-a}}{(1-a)(2-a)}\right]^k.
  \]
  Its sum is $\exp(2t^{3-a}/((1-a)(2-a)))$, whose time exponent
  $3-a$ exceeds the target $(3-a)/2$.
\end{nogo}
\begin{proof}
  The spatial replacement gives $(2t)^k/2^{2k}=t^k/2^k$ for each pairing.
  After this replacement the sum over all pairings is symmetric in the
  $2k$ times. Its integral over a split $(m,n)$ simplex equals the cube
  integral divided by $m!n!$. The identities
  \[
    \sum_{m+n=2k}\frac1{m!n!}=\frac{2^{2k}}{(2k)!},\qquad
    P_k=\frac{(2k)!}{2^k k!},\qquad
    \int_0^t\int_0^t|s-r|^{-a}\,\ud s\ud r
    =\frac{2t^{2-a}}{(1-a)(2-a)}
  \]
  give the stated product and its exponential sum.
  This proves No-go~\ref{P:nogo-06}.
\end{proof}

This is a valid upper bound and a failure of its scale, not a moment lower
bound. Symmetry itself remains useful. The reopening condition is a spatial
estimate with enough order dependence retained before temporal
symmetrization; Proposition~\ref{P:d1-upper} retains this order dependence
and sums the temporal pairings afterwards (see
Remark~\ref{R:d1-upper-checks}).

\subsection{NG-07: block labels are not chronological ranks}
\label{S:nogo-07}

Temporal singularities favor close paired times. A proposed penalty for
long pairs must therefore measure time separation or actual chronological
rank. Labeling all vertices of one chain before those of the other creates
an artificial distance unrelated to either.

\setcounter{nogo}{6}
\begin{nogo}[Large label span does not imply a cubic denominator]
  \label{P:nogo-07}
  In white space with $0<a<1$, consider two length-$k$ chains and the matching
  $j\leftrightarrow k+j$. Although every pair has label distance $k$, its
  unit-time diagram satisfies
  \[
    J_k(1)\ge \left(\frac{3\,4^a}{128}\right)^k k^{-(3-a)k}.
  \]
  Thus no estimate assigning $C^k/(k!)^3$ to these diagrams solely on account
  of their block-label span can hold uniformly in $k$.
\end{nogo}
\begin{proof}
  Let the chain times be $s_j,r_j$, with increments $h_j,g_j$ from the
  common root. White contractions identify their paired positions, and
  integration of the successive spatial increments gives
  $2^{-k}\prod_{j=1}^k\min(h_j,g_j)$. Restrict both $s_j,r_j$ to
  $((j-1/4)/k,j/k)$. The domain has volume $(4k)^{-2k}$,
  each increment is at least $3/(4k)$, and each paired difference is at most
  $1/(4k)$. Multiplication gives the lower bound. The contraction can again
  be obtained at fixed $k$ by Gaussian smoothing; the spatial box bound and
  integrable disjoint temporal factors justify the limit.
  Stirling's formula implies
  $\log((k!)^3J_k(1))\ge a k\log k+O(k)$, contradicting any fixed $C$.
  This proves No-go~\ref{P:nogo-07}.
\end{proof}

% [held back: collaborators' unpublished work]

\subsection[NG-10: crossing covariance survives merging]{NG-10: an unweighted merge does not remove a crossing covariance}
\label{S:nogo-10}

A two-wave convolution estimate can remove an interior chain vertex when
no other factor depends on it. In a crossing pairing that vertex still
interacts with a remote partner. The following order-four calculation
identifies the missing weight before any all-order iteration is attempted.

\setcounter{nogo}{9}
\begin{nogo}[No finite spatial-supremum pullout in the merge lens]
  \label{P:nogo-10}
  In $d=2$, take $\Gamma(z)=c|z|^{-\alpha}$, $c>0$,
  $0<\alpha_0<1/2$ and $3/2\le\alpha<2-\alpha_0$.
  For the root-forward $(0,4)$ diagram
  with pairs $(1,3),(2,4)$, merging vertex two requires the weighted integral
  \[
    W=c\int_a^b(T-v)^{-\alpha_0}
    \int_{\mathbb R^2}G_{v-a}(y-x)G_{b-v}(z-y)|w-y|^{-\alpha}
    \,\ud y\ud v,
    \qquad 0<a<b<T.
  \]
  Here the letters $a,b,T$ are vertex times, so the temporal exponent is
  written $\alpha_0$. The spatial weight has no finite uniform essential
  supremum on admissible merge lenses. Hence, pulling it out by a uniform
  supremum cannot reduce $W$ to an unweighted two-wave convolution.
\end{nogo}
\begin{proof}
  Set $t=5$, $a=1$, $b=3$, $T=4$, $x=(0,0)$,
  $z=(1/2,0)$ and $w=(1/4,0)$. For $v$ near $2$, the point $y=w$
  is strictly inside both disks $|y-x|<v-a$ and $|z-y|<b-v$.
  It also satisfies the outer constraint $|w-z|<T-b$, and $|x|<a$.
  There is therefore an interior ball about $w$ in the merge lens.
  For every $M>0$ a smaller ball has positive Lebesgue measure and
  $|w-y|^{-\alpha}>M$ there. This proves infinite essential supremum;
  the strict inequalities persist under small changes of the fixed data.
  The other pair contributes $(b-a)^{-\alpha_0}c|z-x|^{-\alpha}$ outside
  the merge, whereas the two factors for pair $(2,4)$ depend on $v,y$
  and must remain in $W$. The temporal weight alone admits the bound
  $(T-v)^{-\alpha_0}\le(T-b)^{-\alpha_0}$.
  This proves No-go~\ref{P:nogo-10}.
\end{proof}

The unweighted estimate of~\cite{bolanos-guerrero.nualart.ea:21:averaging}
cannot be substituted for this different object without controlling the
weight. Infinite essential supremum does not imply divergence of $W$ or
of the diagram: the Riesz singularity is locally integrable for $\alpha<2$.
A reopening condition is a covariance-weighted convolution inequality,
or an explicit envelope integrable against the remaining diagram factors,
with the remote-gap cost retained. No such inequality is asserted here.

\subsection{NG-17: direct copying of the first-moment clock undercounts}
\label{S:nogo-17}

The source's first-moment clock representation suggests attaching one
Brownian clock to each ordered chain. A degree-two cross pair already tests
whether this copying preserves the interaction. It does not, even when
the two clocks are shared.

\setcounter{nogo}{16}
\begin{nogo}[Constant-correlation clock copying misses a positive factor]
  \label{P:nogo-17}
  Under the nonzero covariance hypotheses of
  Proposition~\ref{P:clock-copy-obstruction}, let $K_\Gamma(\theta)$ be the
  true cross-pair transform and $J_{\Gamma,\rho}(\theta)$ the direct
  unweighted candidate~\eqref{E:naive-clock-candidate}, $-1\le\rho\le1$.
  Then for every $\theta>0$,
  \[
    0<J_{\Gamma,\rho}(\theta)\le J_{\Gamma,1}(\theta)<K_\Gamma(\theta).
  \]
  If $\alpha=1$, then $J_{\Gamma,1}=(2-2^a)K_\Gamma$.
\end{nogo}
\begin{proof}
  Recall $\mathcal R_\Gamma(q)=\int(q^2+|\xi|^2)^{-1}\mu(\ud\xi)$.
  The exact transform identities proved in
  Proposition~\ref{P:clock-copy-obstruction} give
  \[
    w_\theta(\lambda)=
    \frac{\lambda^{a-2}}{2\theta\Gamma_{\mathrm E}(a)}
    [\mathcal R_\Gamma(\theta)-\mathcal R_\Gamma(\theta+\lambda)],
    \qquad
    K_\Gamma=\int_0^\infty w_\theta(\lambda)\,\ud\lambda,
  \]
  \[
    J_{\Gamma,1}=\int_0^\infty
    \frac{2\theta}{2\theta+\lambda}w_\theta(\lambda)\,\ud\lambda.
  \]
  The density is strictly positive because $\mu\ne0$ and the resolvent
  strictly decreases. It is integrable: homogeneity bounds it by multiples
  of $\lambda^{a-1}$ near zero and $\lambda^{a-2}$ at infinity.
  The multiplier is strictly between zero and one. For smaller correlations,
  $u+v-2\rho\min(u,v)\ge|u-v|$, and the positive Laplace integral
  $\Psi_a(r,v)$ in~\eqref{E:averaged-clock-kernel} decreases with $v$.
  This proves the chain of inequalities. When $\alpha=1$, insert
  $\mathcal R_\Gamma(q)=\mathcal R_\Gamma(1)/q$ and use
  $\int_0^\infty x^{a-1}/(x+c)\,\ud x
  =c^{a-1}\Gamma_{\mathrm E}(a)\Gamma_{\mathrm E}(1-a)$ to obtain the ratio.
  This proves No-go~\ref{P:nogo-17}.
\end{proof}

This is the already proved book obstruction, not a new clock impossibility
theorem. At the static endpoint the discrepancy disappears, so homogeneity
and static checks alone cannot detect it. At degree two the missing
shared-clock multiplier is exactly $1+\lambda/(2\theta)$; no higher-order
correction rule follows. Reopening requires relative-order weights or a
different coupling or representation that reproduces the actual cross-pair
transform. Remark~\ref{R:clock-copy-scope} gives the full boundary of this
statement, including the distinction between gap transforms and physical
unequal-time moments.

\subsection[NG-18: an oversized middle-cross envelope]{NG-18: erasing the middle-cross shift makes an envelope too large}
\label{S:nogo-18}

A later white-space comparison retains four vertices and a middle cross
pair. Removing the clock-overlap penalty gives a positive upper bound.
One further simplification, erasing the shift of its root kernel, also
increases the bound, but increases it beyond the comparator at an explicit
parameter point. This is an obstruction to closing that envelope, not a
counterexample to the original comparison.

\setcounter{nogo}{17}
\begin{nogo}[The shift-erased envelope is not below the wave transform]
  \label{P:nogo-18}
  Consider the white $(1,3)$ diagram with middle-cross partition $02\cdot13$,
  common gap damping one, cross rate $\lambda\ge0$ and self rate $\nu\ge0$.
  Put $b=1+\nu$, $c=1+\lambda$, $A=\lambda+\lambda^2/2$. Its actual wave
  transform $I$ and the shift-erased clock upper envelope $M_b(0;A)$ satisfy
  at $b=1,c=11/10$
  \[
    I=\frac{545}{11616},\qquad M_b(0;A)=\frac{1025}{21483},\qquad
    I-M_b(0;A)=-\frac{6005}{7562016}<0.
  \]
  Consequently the proposed pointwise closure $M_b(0;A)\le I$ is false.
\end{nogo}
\begin{proof}
  We give the integral identifications so that the rational test specifies
  which bound fails. With singleton gap $u$ and long-chain gaps $v,w,z$,
  the spatial factor is
  $\int G_u(X)G_v(Y)G_w(X-Y)G_z(Y-X)\,\ud X\ud Y$.
  Set $R=u+v$, $q=u-v$. Integration in $R$, including its Jacobian,
  gives the exact wave expression
  \[
    I=\frac1{16}\int_{w,z>0}\int_{\mathbb R}
    e^{-b(w+z)-\lambda|q-w|}
    \int_0^{\min(w,z)}e^{-\max(|q|,t)}\,\ud t\ud q\ud w\ud z.
  \]
  To check the spatial reduction, the derivative of the area
  $\{|X|<u,|Y|<v,|X-Y|<C\}$ with respect to $C$ is
  $2(R-\max(C,|q|))_+$; the area vanishes at $C=0$.
  Integrate this against $e^{-R}/32$ for $R>|q|$.
  Next integrate $z\ge t$ and put $w=t+x$. Splitting the remaining $q$
  integral into $q<0$, $0\le q\le t$, and $q\ge t$ gives respectively
  \[
    \frac{1/c+1/(2b+\lambda+1)}{16b(b+\lambda)(2b+2\lambda+1)},\quad
    \frac1{16b(b+\lambda)(2b+1)(2b+1+\lambda)},
  \]
  \[
    \frac{1/(b+\lambda)+1/c}{16b(2b+1)(b+1)}.
  \]
  Adding gives
  \[
    I=\frac{2b^2+4bc+2c^2-c+1}
    {8bc(b+1)(2b+1)(b+c-1)(2b+2c-1)}.
  \]

  For the clock envelope write $h(r)=(2\pi r)^{-1/2}$ and
  $m(s)=1/(\pi\sqrt2\sqrt{s-1/2})$, $s>1/2$.
  The two spatial pair differences have covariance matrix
  $\left(\begin{smallmatrix}u+v+w&-w\\-w&w+z\end{smallmatrix}\right)$.
  Writing $r=w+z$, its determinant is $rR+wz$, so their joint white
  density factors as $h(r)h(R+wz/r)$. After deleting the nonnegative
  clock-overlap penalty, the root kernel is
  \[
    H_\gamma(q)=\int_{|q|}^\infty e^{-R/2}h(R+\gamma)\,\ud R
    =\int_{1/2}^\infty\frac{m(s)}s
    e^{-(s-1/2)\gamma}e^{-s|q|}\,\ud s,
    \qquad \gamma=\frac{wz}{r}.
  \]
  The equality follows from the positive Gamma integral for $h$.
  Its convolution
  with $e^{-A|q-w|}$ equals $2S_{A,s}(w)/(s+A)$, where
  \[
    S_{A,s}(w)=e^{-Aw}+\int_0^w A e^{-At}e^{-s(w-t)}\,\ud t\le1.
  \]
  Replacing $S$ by one is precisely the shift erasure. More generally,
  replacing it by $e^{-\eta w}$ defines
  \[
    M_b(\eta;A)=\frac1{16}\int_{1/2}^\infty
    \frac{m(s)}{s(s+A)}\int_0^\infty h(r)e^{-b^2r/2}
    \int_0^r e^{-(s-1/2)w(r-w)/r-\eta w}\,\ud w\ud r\ud s.
  \]
  All integrals here are positive. Put $w=rt$ and $p=\sqrt{b^2+2\eta}$.
  The elementary bridge integral is
  \[
    \int_0^1[b^2+2\eta t+2\tau t(1-t)]^{-3/2}\,\ud t
    =\frac{2(p+b)}{bp((p+b)^2+2\tau)}.
  \]
  It follows by integrating a quadratic power; at $\eta=\tau=0$ both sides
  are $b^{-3}$. Substitution $z=\sqrt{2s-1}$ and partial fractions give
  \begin{align*}
    M_b(\eta;A)&=\frac{p+b}{2\pi bp}
    \int_0^\infty\frac{\ud z}
    {(z^2+1)(z^2+c^2)(z^2+(p+b)^2)}\\
    &=\frac{1+c+p+b}{4bp c(c+1)(c+p+b)(1+p+b)}.
  \end{align*}
  The usual identity $\int_0^\infty(z^2+q^2)^{-1}\ud z=\pi/(2q)$
  proves the partial fraction evaluation; coincident rates follow by
  continuity. At $\eta=0$, $p=b$. Substituting the asserted rational
  parameters in these two expressions gives the displayed strict gap.
  This proves No-go~\ref{P:nogo-18}.
\end{proof}

The boundary value $\nu=0$ is a legitimate fixed-parameter transform;
continuity of both rational expressions also gives an overshoot for some
$\nu>0$ nearby. The no-go does not refute $J\le I$, where $J$ is the
original shared-clock candidate, nor does it decide an integrated
comparison using this envelope. The surviving middle-cross comparison
retains $S_{A,1/2}(w)$, instead of replacing it by one, and obtains a smaller
envelope. Reopening this particular closure requires keeping a shift cost
or proving a different integrated estimate. The historical standalone
record was not separately adjudicated; the subsequent accepted
middle-cross analysis and this reconstructed proof do not change that
historical review status.

\section{Covariance transformations and their limits}\label{S:nogo-covariance}

Covariance comparisons offer a way to use a simpler model without estimating
all the pairings again. Three different requirements enter such an argument:
positive definiteness, local finiteness, and uniformity of the comparison.
None follows from homogeneity alone. The examples below separate these
requirements. The final example also distinguishes positive semidefinite
order from pointwise order, which matters for the wave mean.

\subsection{NG-08: the translated spatial maximum}\label{S:nogo-08}

The conditional higher-dimensional argument replaces a translated wave
convolution by a uniform spatial maximum. This would remove the dependence
on the other vertices of a pairing. The price is an additional hypothesis,
not a consequence of the covariance assumptions.

\setcounter{nogo}{7}
\begin{nogo}[A spatial maximum need not be finite]\label{P:nogo-08}
  Hypothesis~\ref{H:spatial-maximum} fails for two-dimensional Riesz
  covariance of degree $-\alpha$ when $3/2\le\alpha<2$. For each
  $1<\alpha<2$, it also fails for some nonnegative, positive definite,
  locally integrable homogeneous product covariance in dimensions two
  and three. These examples can satisfy the model condition
  $\alpha+\alpha_0<2$.
\end{nogo}
\begin{proof}
  We exhibit the boundary divergence for planar Riesz covariance and then the
  anisotropic product example in both dimensions. The smoothing limit shows
  why neither example admits the proposed uniform spatial maximum.


  \medskip
  \noindent\textit{Step 1.~}We test the translated planar Riesz convolution.
  We restate the witnesses from Proposition~\ref{P:higher-obstructions}.
  For Riesz covariance translate the singularity to $e_1$, on the boundary
  of the unit wave disk. Write $y=e_1+\rho(\cos\phi,\sin\phi)$ and restrict
  to $0<\rho<1/2$, $2\pi/3\le\phi\le4\pi/3$. Then
  $0<1-|y|^2\le2\rho$, so the wave density and polar Jacobian give
  \[
    \int G_1^{(2)}(\ud y)c|y-e_1|^{-\alpha}
    \ge \frac{c}{3\sqrt2}\int_0^{1/2}\rho^{1/2-\alpha}\ud\rho
    =\infty.
  \]
  At $\alpha=3/2$ the divergence is logarithmic.


  \medskip
  \noindent\textit{Step 2.~}We test the product covariance and its smoothed
  limits.
  For the product example choose
  \[
    a_1=\frac{1+\max\{2-\alpha,\alpha-1\}}2,
    \qquad a_2=\alpha-a_1.
  \]
  Then $0<a_1,a_2<1$ and $a_1+a_2/2>1$. The function
  $\gamma(x)=c|x_1|^{-a_1}|x_2|^{-a_2}$, constant in $x_3$ when $d=3$,
  is locally integrable and homogeneous. Each factor is a positive
  Gaussian mixture, hence the product is positive definite. With
  $u=1-y_1$ and $y_2=\sqrt{2u-u^2}\,w$, direct integration yields
  \[
    \int G_1^{(2)}(\ud y)\gamma(y-e_1)
    =\frac{c B((1-a_2)/2,1/2)}{2\pi}
    \int_0^2u^{-a_1}(2u-u^2)^{-a_2/2}\ud u=\infty.
  \]
  The projection of the three-dimensional wave measure onto its first
  two coordinates is $G_1^{(2)}$, which proves the three-dimensional claim.
  Gaussian smoothing converges off the singular sets, which have zero
  translated wave measure. Fatou's lemma therefore gives divergence in
  lower limit as the smoothing widths decrease to zero. Thus, smoothing
  cannot restore the required uniform constant. Taking
  $0<\alpha_0<2-\alpha$ ensures admissibility.
  This proves No-go~\ref{P:nogo-08}.
\end{proof}

This result excludes the spatial-supremum input, not moment finiteness or
the required time scale. The conditional theorem remains useful where its
hypothesis holds. Reopening this route for the full class would require
an integrated estimate retaining the geometry lost by the supremum;
Section~\ref{S:interval-upper} supplies a different full-class argument.

\subsection{NG-09: radial lowering in three dimensions}\label{S:nogo-09}

Multiplication by $|x|^{\alpha-1}$ gives the desired endpoint degree $-1$.
The following example shows why this degree calculation does not construct
a covariance. The original endpoint review returned MEDIUM confidence and
was not accepted at the HIGH threshold. The argument below is a complete
reconstruction; it does not change that historical disposition.

\setcounter{nogo}{8}
\begin{nogo}[Radial lowering can destroy positive definiteness]\label{P:nogo-09}
  In dimension three let
  \[
    \Gamma(x)=|x_1|^{-9/10}|x_2|^{-3/5},\qquad
    F(x)=|x|^{1/2}\Gamma(x).
  \]
  The covariance $\Gamma$ satisfies the homogeneous model assumptions with
  $\alpha=3/2$ and $\alpha_0=1/4$. The locally integrable tempered function
  $F$ is not positive definite as a distribution.
\end{nogo}
\begin{proof}
  We verify the original covariance, then show that the radially modified
  function has a heat convolution growing along the third coordinate. This
  contradicts the boundedness that positive definiteness would require.


  \medskip
  \noindent\textit{Step 1.~}We check the original covariance and the modified
  function.
  The exponents of both coordinate factors lie in $(0,1)$. Local
  integrability, temperedness, and positive definiteness of $\Gamma$
  follow from the one-dimensional Gaussian mixture formula, applied in
  each coordinate. The third factor is the constant covariance one.
  Multiplication by $|x|^{1/2}$ preserves local integrability and temperedness.


  \medskip
  \noindent\textit{Step 2.~}We obtain the contradiction after Gaussian
  convolution.
  Fix a centered Gaussian heat density $p_\sigma$. Its convolution with
  $F$ is finite at zero: the coordinate singularities are integrable and
  the remaining growth is polynomial. On the other hand, restrict the
  convolution at $Le_3$ to
  \[
    1\le y_1\le2,\qquad 1\le y_2\le2,\qquad |y_3-L|\le1.
  \]
  For $L\ge2$, the Gaussian factor $p_\sigma(Le_3-y)$ is bounded below
  there by a fixed positive constant, the coordinate powers are bounded
  below, and $|y|^{1/2}\ge(L-1)^{1/2}$. Hence,
  \[
    (p_\sigma*F)(Le_3)\ge c_\sigma\sqrt L.
  \]
  If $F$ were positive definite, its quadratic form on Schwartz tests
  would satisfy Cauchy--Schwarz. Take a Gaussian test whose autocorrelation
  is $p_\sigma$, and its translate by $x$. The two diagonal quadratic
  forms are both $(p_\sigma*F)(0)$, while the cross form is
  $(p_\sigma*F)(x)$. Consequently,
  \[
    |(p_\sigma*F)(x)|\le(p_\sigma*F)(0),\qquad x\in\mathbb R^3.
  \]
  The lower bound along $Le_3$ contradicts this inequality.
  This proves No-go~\ref{P:nogo-09}.
\end{proof}

The obstruction concerns this Euclidean radial multiplication. It leaves
coordinatewise constructions and the planar endpoint argument intact;
see Remark~\ref{R:three-endpoint-scope}. A successor must verify positive
definiteness of its endpoint, rather than infer it from degree and sign.

\subsection{NG-11: additive variance interpolation}\label{S:nogo-11}

A covariance mixture would be easy to handle if its second moment were
log-convex in the additive variance. The following scalar calculation
refutes that unrestricted principle. It is outside the strict homogeneous
model and is not a counterexample to a model-specific comparison.

\setcounter{nogo}{10}
\begin{nogo}[A scalar wave moment is not log-convex in variance]\label{P:nogo-11}
  Let $y''=Zy$, $y(0)=1$, $y'(0)=0$, where $Z$ is centered Gaussian of
  variance $A$. For $F(A)=\E[y(1)^2]$,
  \[
    F(100)^2>F(99)F(101).
  \]
  Thus, $\log F$ is not convex on $(0,\infty)$.
\end{nogo}
\begin{proof}
  We express the scalar moment as a positive series and bound its tail. The
  resulting finite rational enclosures disprove the stated log-convexity
  inequality.


  \medskip
  \noindent\textit{Step 1.~}We obtain the series and a controlled tail.
  Interpreting $\cosh\sqrt z$ by its entire power series gives
  \[
    F(A)=1+\sum_{k\ge1}t_k(A),\qquad
    t_k(A)=\frac{2^{3k-1}(2k)!}{k!(4k)!}A^k.
  \]
  Absolute termwise Gaussian integration follows from
  $\E\exp(2\sqrt{|Z|})<\infty$. All displayed terms are positive.
  Their ratio satisfies
  \[
    \frac{t_{k+1}(A)}{t_k(A)}
    =\frac{16A(2k+1)}{(4k+1)(4k+2)(4k+3)(4k+4)}
    \le\frac{16A}{(4k+1)^3}.
  \]
  Write $S_{20}=1+\sum_{k=1}^{20}t_k(A)$ and
  $q_A=16A/81^3<1$ for the three values in question. Then
  \[
    S_{20}\le F(A)\le S_{20}+\frac{t_{21}(A)}{1-q_A}.
  \]

  \medskip
  \noindent\textit{Step 2.~}We compare the rigorous rational enclosures.
  Substitution in these finite rational expressions gives the following
  outward enclosures, with common denominator $10^6$:
  \[
    \begin{array}{c|rr}
      A&10^6\text{ times lower bound}&10^6\text{ times upper bound}\\ \hline
      99&207442589&207442590\\
      100&212332413&212332414\\
      101&217304018&217304019
    \end{array}
  \]
  In particular,
  \[
    F(100)^2-F(99)F(101)
    \ge\frac{6945091633359}{10^{12}}>0.
  \]
  This proves No-go~\ref{P:nogo-11}.
\end{proof}

Positive power-series coefficients do imply convexity of
$t\mapsto\log F(e^t)$, by H\"older's inequality. Thus, geometric
interpolation of amplitudes survives. An additive-covariance argument
would need additional hypotheses and a separate proof; the scalar test
prevents using log-convexity as a general Gaussian-wave fact.

\subsection{NG-12: a heat--Mellin square need not be locally finite}
\label{S:nogo-12}

Squaring a heat-smoothed covariance preserves positive definiteness, and
a Mellin weight can prescribe homogeneity. This suggests an endpoint
recipe with no angular decomposition. Its missing requirement is local
finiteness of the infinite positive mixture.

\setcounter{nogo}{11}
\begin{nogo}[Failure of a heat--Mellin endpoint recipe]\label{P:nogo-12}
  Let $d=3$, $\alpha=11/10$, $\alpha_0=1/4$, and
  \[
    \Gamma(x)=|x_1|^{-9/10}|x_2|^{-1/5},\qquad F_s=p_s*\Gamma.
  \]
  Although $\Gamma$ is an admissible covariance, the putative endpoint
  density of degree $-\beta$, with $\beta=9/5$,
  \[
    H(x)=\int_0^\infty
    s^{(2\alpha-\beta)/2-1}F_s(x)^2\ud s
    =\int_0^\infty s^{-4/5}F_s(x)^2\ud s
  \]
  does not define a locally finite measure.
\end{nogo}
\begin{proof}
  The coordinate exponents show admissibility as in No-go~\ref{P:nogo-09}.
  Heat convolution factorizes in the first two coordinates. For
  $0<s<1/16$ and $|x_1|\le\sqrt s$, integrating the first Gaussian
  convolution over $\sqrt s\le y_1\le2\sqrt s$ gives
  \[
    (p_s^{(1)}*|\,\cdot\,|^{-9/10})(x_1)\ge c s^{-9/20}.
  \]
  For $1/2\le x_2\le1$, integrate the second convolution over
  $x_2\le y_2\le x_2+\sqrt s$; its value is bounded below by a positive
  constant. The third coordinate contributes one. On the fixed box
  $B=[-1,1]\times[1/2,1]\times[-1,1]$, the strip
  $|x_1|\le\sqrt s$ therefore gives
  \[
    \int_B F_s(x)^2\ud x\ge c s^{1/2-9/10}=c s^{-2/5}.
  \]
  Tonelli's theorem implies
  \[
    \int_B H(x)\ud x\ge c\int_0^{1/16}s^{-6/5}\ud s=\infty.
  \]
  This proves No-go~\ref{P:nogo-12}.
\end{proof}

The positive-definiteness calculation for finite mixtures is not the
problem. The infinite mixture has infinite mass on a compact set. This
excludes the displayed recipe, not endpoint interpolation itself. A
replacement must prove local finiteness, for example by controlling the
coordinate singularities before integrating the smoothing scale.

\subsection{NG-13: angular smoothing and isotropic envelopes}\label{S:nogo-13}

Smooth angular densities admit finite angular suprema. It is tempting to
approximate a singular covariance by such densities and apply the Riesz
estimate with that supremum as amplitude. The limit below shows why the
amplitudes cannot remain bounded.

\setcounter{nogo}{12}
\begin{nogo}[No uniform isotropic envelope for a plane covariance]
  \label{P:nogo-13}
  In $\mathbb R^3=\mathbb R^2\times\mathbb R$ set
  \[
    \Gamma(\ud y,\ud z)=\frac1{2\pi}|y|^{-1/2}\ud y\,\delta_0(\ud z).
  \]
  This is an admissible covariance with $\alpha=3/2$ and
  $\alpha_0=1/4$. If nonnegative measures $\Gamma_n$ converge vaguely to
  $\Gamma$ and satisfy
  \[
    \Gamma_n(\ud x)\le C_n|x|^{-3/2}\ud x,
  \]
  then $C_n\to\infty$.
\end{nogo}
\begin{proof}
  We verify that the plane-supported covariance is admissible. A bounded
  subsequence of isotropic envelope constants would force its vague limit to
  have a Lebesgue density, contradicting its planar mass.


  \medskip
  \noindent\textit{Step 1.~}We verify the singular covariance example.
  The planar Riesz factor is locally integrable and positive definite by
  its Gaussian mixture representation. The normal Dirac covariance is
  positive definite: in its quadratic form the two normal coordinates
  coincide, leaving an integral of nonnegative tangential quadratic
  forms. Thus, the product is positive definite. For every $r>0$ and
  Borel set $A$, the substitution $y=ry'$ in its defining integral gives
  $\Gamma(rA)=r^{2-1/2}\Gamma(A)=r^{3/2}\Gamma(A)$. Its finite mass on
  a radius $R$ ball is proportional to $R^{3/2}$; polynomial mass growth
  also gives temperedness.


  \medskip
  \noindent\textit{Step 2.~}We contradict bounded envelope constants by
  singular support.
  Suppose a subsequence of $C_n$ were bounded by $C<\infty$. For every
  nonnegative compactly supported continuous $f$, vague convergence gives
  \[
    \int f\ud\Gamma=\lim_n\int f\ud\Gamma_n
    \le C\int f(x)|x|^{-3/2}\ud x.
  \]
  This is measure domination, so $\Gamma$ would be absolutely continuous
  with respect to three-dimensional Lebesgue measure. But the compact
  planar annulus $\{(y,0):1\le|y|\le2\}$ has positive $\Gamma$ mass and
  zero Lebesgue measure. This contradiction excludes every bounded
  subsequence, proving the asserted limit.
  This proves No-go~\ref{P:nogo-13}.
\end{proof}

This argument applies in particular to angular smoothing whenever the
smoothed covariances converge to the plane measure. It requires no
ellipsoid representation. It rules out a uniform scalar isotropic density
envelope, not other covariance majorants or density of any proposed cone.
The reopening condition is an estimate that remains uniform without taking
an angular supremum, or a comparison adapted to the singular support.

\subsection{NG-19: discarding an independent Gaussian residual}
\label{S:nogo-19}

An orthogonal Gaussian projection produces an independent residual with
positive semidefinite covariance. This makes projection a useful change-of-law
tool. It does not give pointwise nonnegative residual covariance, and the
wave mean need not increase when that residual is restored.

\setcounter{nogo}{18}
\begin{nogo}[A positive semidefinite residual can lower the wave mean]
  \label{P:nogo-19}
  There is an admissible one-dimensional parent covariance $C$ and an
  actual rank-one Gaussian projection with covariance $P$ such that
  $C-P$ is positive semidefinite but
  \[
    \E u_C(h;1,0)<\E u_P(h;1,0)
  \]
  for all sufficiently small positive noise couplings $h$. The projection
  $P$ is not asserted to be stationary or homogeneous.
\end{nogo}
\begin{proof}
  We construct an actual Gaussian projection of the parent field and compare
  the two degree-two mean coefficients. Bounds on the higher-order tails then
  turn that coefficient gap into a strict small-coupling failure of residual
  monotonicity.


  \medskip
  \noindent\textit{Step 1.~}We construct the parent field and its rank-one
  projection.
  Put $a=99/100$ and
  \[
    C((s,x),(t,y))=(1-a)|s-t|^{-a}|x-y|^{-1/2}.
  \]
  Its spatial degree is $\alpha=1/2$, and $\alpha+a<2$. Gaussian mixtures
  of the two power kernels verify positive definiteness and local
  integrability. Let $\psi(s,x)=\mathbf1_{[0,1]}(s)\mathbf1_{[-1,1]}(x)$,
  and $V=\E W(\psi)^2$. The indicator is a legal Hilbert-space test by
  bounded smooth approximation and dominated convergence. Local integrability
  and positivity of $C$ give $0<V<\infty$. Put $Z=W(\psi)/\sqrt V$ and
  define the covariance profile by the locally integrable kernel formula
  \[
    m(s,x)=\frac1{\sqrt V}\int_0^1\int_{-1}^1
    C((s,x),(r,y))\ud y\ud r.
  \]
  For every smooth compactly supported spacetime test $f$, Fubini's
  theorem identifies $\langle m,f\rangle=\E[W(f)Z]$.
  Thus, $m$ is the covariance of the generalized field with $Z$; no
  pointwise random variable $W(s,x)$ is used. Gaussian projection gives
  $W=Zm+W^\perp$ with $W^\perp$ independent of $Z$ and
  covariance $C-P$, where $P=m\otimes m$. Cauchy--Schwarz gives its
  positive semidefiniteness.


  \medskip
  \noindent\textit{Step 2.~}We bound the projected profile and compare the
  mean coefficients.
  Direct integration gives, on $0\le s\le1$,
  \[
    E_\phi=\frac{16\sqrt2}{3},\quad V=\frac{2E_\phi}{2-a},\quad
    m(s,x)=\frac{B_a(s)F(x)}{\sqrt V},
  \]
  where $B_a(s)=s^{1-a}+(1-s)^{1-a}\ge1$ and
  $F(x)=\int_{-1}^1|x-y|^{-1/2}\ud y$. We have $F\le4$ everywhere and
  $F(x)\ge2\sqrt2$ for $|x|\le1$. Thus, $m$ is bounded and nonnegative,
  and $m(s,x)m(t,y)\ge8/V$ on the backward cone of $(1,0)$.

  For a covariance $D$ the degree-two ordinary mean coefficient is
  \[
    A_D=\int_{0<s<t<1}\int_{\mathbb R^2}
    G_{1-t}(-y)G_{t-s}(y-z)D((s,z),(t,y))
    \ud z\ud y\ud s\ud t.
  \]
  This is the single self-pair term of Proposition~\ref{P:mixed}; there
  is no additional factorial. With $G_r(x)=\mathbf1_{\{|x|<r\}}/2$,
  integration of the parent covariance yields
  \[
    A_C=\frac{1-a}{1-\alpha}
    \frac1{(2-\alpha-a)(3-\alpha-a)(4-\alpha-a)}
    =\frac{20000}{1932951}<\frac1{96}.
  \]
  For covariance one the same integral is
  $\int_{0<s<t<1}(1-t)(t-s)\ud s\ud t=1/24$. Therefore,
  \[
    A_P\ge\frac{8}{24V}
    =\frac{101}{3200\sqrt2}>\frac1{48}.
  \]


  \medskip
  \noindent\textit{Step 3.~}We pass the strict coefficient gap to the full
  means.
  We justify passage from coefficients to full means. Write the parent
  solution as $u_C(h)=\sum_{n\ge0}h^nX_n$. For any fixed $r>1$,
  finiteness of the second moment at amplified covariance $r^2C$ and
  nonnegativity of the mixed moments in Proposition~\ref{P:mixed} give
  \[
    \|X_n\|_2\le \sqrt{B_r}\,r^{-n},\qquad
    B_r=\E|u_C(r;1,0)|^2<\infty.
  \]
  Hence, the $L^2$ norm of the tail starting at $n=4$ is at most
  $\sqrt{B_r}(|h|/r)^4/(1-|h|/r)$. Odd means vanish by Gaussian sign
  symmetry, giving $\E u_C(h)=1+h^2A_C+O(h^4)$.

  For $P$, the potential is the bounded profile $Zm$. If $M=\sup|m|$,
  the deterministic Picard coefficients at time one are bounded by
  $M^n/(2n)!$. The absolute mean series for $|h|\le1$ is bounded by
  $\E\cosh\sqrt{M|Z|}<\infty$. Factoring $h^4$ from its tail and
  using symmetry therefore gives $\E u_P(h)=1+h^2A_P+O(h^4)$.
  The strict coefficient gap proves the claimed inequality.
  This proves No-go~\ref{P:nogo-19}.
\end{proof}

The same distinction persists under a variance tilt of the projected
coordinate. For $1<q<2$, its density relative to the parent law is
\[
  L_q(Z)=q^{-1/2}\exp\bigl((1-q^{-1})Z^2/2\bigr),\qquad
  \E L_q^2=[q(2-q)]^{-1/2}.
\]
The equivalent Gaussian law has covariance $Q=C+(q-1)P$. The same original
measurable solution has an $O(h^4)$ mean remainder under this law by
Cauchy--Schwarz and the preceding $L^2$ tail estimate. Its degree-two
coefficient is $A_Q=A_C+(q-1)A_P<qA_P-1/96$. This is an equivalent-law
statement, not a new stationary-model existence theorem for $Q$.
To identify its degree-two coefficient, one can use the same regularization
as the parent construction. If $K_\delta$ denotes the truncated temporal
covariance and $p_\epsilon$ the spatial heat kernel, its cross-covariance
with $Z$ is
\[
  m_{\epsilon,\delta}(s,x)
  =\frac{(1-a)\int_0^1K_\delta(s-r)\ud r}{\sqrt V}
  (p_\epsilon*F)(x).
\]
The bounds $0\le K_\delta\le|\,\cdot\,|^{-a}$ and $F\le4$ give
$0\le m_{\epsilon,\delta}\le8/\sqrt V$. These profiles converge to
$m$ at interior times. Dominated convergence on the finite-mass terminal
wave kernel identifies the added $P$ contribution. The parent degree-two
term converges in $L^2$, and $L_q\in L^2$ transfers its expectation to the
equivalent law. Thus, the fixed-order Gaussian covariance identity indeed
gives the stated $A_Q$, without invoking a homogeneous representation for
$Q$. Consequently, discarding the residual also fails in this tilted comparison.
A large-time bridge for specially chosen projections could still hold,
but it must be proved at the required scale. Positive semidefinite order
alone cannot supply it.

\section{Phase propagation and Gaussian controls}\label{S:nogo-phase}

A wave equation transports displacement and velocity together. Its scalar
solution value is not a state variable for composition. Likewise, a
Gaussian change of law transports an entire random solution, not merely
the solution evaluated at the mean of the new noise. These distinctions
explain the failures in this section. The scalar examples are diagnostics
for generic proof principles; their scope is stated separately from the
homogeneous stochastic wave model.

\subsection{NG-14: scalar moments do not compose across blocks}
\label{S:nogo-14}

Independent static time blocks look favorable for multiplying second
moments. Independence does remove correlations between their potentials,
but it does not reset the velocity to zero at a block boundary. The
original record computed formal coefficients only. We include the
integrable-series argument needed to turn that diagnostic into a
small-variance counterexample.

\setcounter{nogo}{13}
\begin{nogo}[Two independent wave blocks do not multiply scalar moments]
  \label{P:nogo-14}
  Let $Z_1,Z_2$ be independent centered Gaussians of variance $A$. Solve
  $y''=Z_1y$ on $[0,1]$ and $y''=Z_2y$ on $[1,2]$, with continuous phase
  at time one and initial phase $(y(0),y'(0))=(1,0)$. Then, as
  $A\downarrow0$,
  \[
    \E y(1)^2=1+\frac A3+O(A^2),\qquad
    \E y(2)^2=1+3A+O(A^2).
  \]
  In particular, $\E y(2)^2>(\E y(1)^2)^2$ for all sufficiently small
  positive $A$.
\end{nogo}
\begin{proof}
  Define the entire functions
  \[
    c(z)=\sum_{n\ge0}\frac{z^n}{(2n)!},\qquad
    s(z)=\sum_{n\ge0}\frac{z^n}{(2n+1)!}.
  \]
  The unit-time phase propagator and the terminal displacement are
  \[
    M(z)=\begin{pmatrix}c(z)&s(z)\\zs(z)&c(z)\end{pmatrix},\qquad
    y(2)=c(Z_2)c(Z_1)+s(Z_2)Z_1s(Z_1).
  \]
  Expanding through total degree two in $(z,w)$ gives
  \[
    c(w)c(z)+s(w)zs(z)
    =1+\frac32z+\frac12w
    +\frac5{24}z^2+\frac5{12}zw+\frac1{24}w^2
    +\text{terms of degree at least three}.
  \]
  After squaring, the coefficients of $z^2$ and $w^2$ are respectively
  $8/3$ and $1/3$. Independence and centering eliminate the $zw$ term.
  Similarly, the coefficient of $z^2$ in $c(z)^2$ is $1/3$.

  For completeness, these expansions admit the stated remainders after
  Gaussian expectation. The series of $c$ and $s$ have nonnegative
  coefficients, and
  \[
    |c(z)|+|s(z)|\le2e^{\sqrt{|z|}}.
  \]
  Consequently, the absolute series for the square of the displayed
  two-block displacement is bounded by
  $C(1+|z|)^2e^{2\sqrt{|z|}+2\sqrt{|w|}}$.
  This is integrable at any fixed Gaussian variance $A_1>0$. Absolute
  integration allows termwise expectation. Odd powers in either variable
  vanish; every remaining term has even total degree and a nonnegative
  coefficient. Thus, the expectation is a convergent series
  $\sum_{j\ge0}b_jA^j$ with $b_j\ge0$ at $A_1$. For $0\le A\le A_1$,
  \[
    \sum_{j\ge2}b_jA^j
    \le (A/A_1)^2\sum_{j\ge2}b_jA_1^j=O(A^2).
  \]
  The same argument applies to the one-block moment. Its square equals
  $1+2A/3+O(A^2)$, whereas the two-block moment equals
  $1+3A+O(A^2)$. Their difference is $7A/3+O(A^2)$.
  This proves No-go~\ref{P:nogo-14}.
\end{proof}

\begin{figure}
  \centering
  \begin{tikzpicture}[x=1cm,y=1cm]
    \node[ham/box,text width=1.8cm] (initial) at (0,0)
    {Initial phase\\$(1,0)$};
    \node[ham/box,text width=3cm] (middle) at (4.1,0)
    {After block 1\\$(c(Z_1),Z_1s(Z_1))$};
    \node[ham/box,text width=2.6cm] (final) at (8.7,0)
    {After block 2\\$M(Z_2)M(Z_1)e_1$};
    \draw[ham/arrow] (initial)--node[above,font=\small] {$M(Z_1)$} (middle);
    \draw[ham/arrow] (middle)--node[above,font=\small] {$M(Z_2)$} (final);
    \node[ham/note,text width=10cm] at (4.4,-1.7)
    {Retained displacement: $c(Z_2)c(Z_1)+s(Z_2)Z_1s(Z_1)$.\\
    Resetting the velocity would retain only $c(Z_2)c(Z_1)$.};
    \node[ham/note,text width=10cm] at (4.4,-3)
    {True second moment: $1+3A+O(A^2)$.\\
    Product of reset-block moments: $1+\frac23A+O(A^2)$.};
  \end{tikzpicture}
  \caption[Velocity survives a block boundary]{Two independent scalar static
    Gaussian blocks, each of variance $A$. The second block acts on both
    components of the first block's phase. This is a scalar diagnostic
    outside the target homogeneous-noise model; see
    No-go~\ref{P:nogo-14}.}
  \label{Fig:nogo-phase-blocks}
\end{figure}


This scalar model has spatially constant noise and independent static
blocks, outside Hypothesis~\ref{H:homog}. It excludes the unrestricted
scalar multiplication rule, not a phase-space operator estimate. Corollary~\ref{C:static-interaction} gives a valid replacement for independent
successive blocks: it controls both phase coordinates and includes explicit
block losses whose dyadic sum is finite. Its stated scope does not identify
products across overlapping levels with the original evolution.

\subsection{NG-15: the missing second-order mean term}\label{S:nogo-15}

Keeping the first random Duhamel term gives a nonnegative squared norm.
It is not the complete coefficient of the variance in the expected phase
energy. The mean of the second Duhamel term contributes at the same order.

\setcounter{nogo}{14}
\begin{nogo}[The first variance coefficient need not be nonnegative]
  \label{P:nogo-15}
  For the scalar phase propagator $M$ in No-go~\ref{P:nogo-14} and a
  standard Gaussian $Z$, write
  \[
    \E\|M(\sigma Z)x\|^2=\|M(0)x\|^2+\sigma^2x^{\mathsf T}Kx
    +O(\sigma^4).
  \]
  Then
  \[
    K=\begin{pmatrix}4/3&4/5\\4/5&17/45\end{pmatrix},\qquad
    (1,-2)K(1,-2)^{\mathsf T}=-\frac{16}{45}.
  \]
  The matrix $K$ is the full first coefficient in the variance parameter,
  not a covariance matrix. Omitting the second-order mean term also
  underestimates this coefficient for $x=(1,0)^{\mathsf T}$.
\end{nogo}
\begin{proof}
  The coefficients of $M(z)=M_0+zM_1+z^2M_2+\cdots$ are
  \[
    M_0=\begin{pmatrix}1&1\\0&1\end{pmatrix},\quad
    M_1=\begin{pmatrix}1/2&1/6\\1&1/2\end{pmatrix},\quad
    M_2=\begin{pmatrix}1/24&1/120\\1/6&1/24\end{pmatrix}.
  \]
  Gaussian expectation of the squared phase norm gives
  \[
    K=M_1^{\mathsf T}M_1+M_0^{\mathsf T}M_2+M_2^{\mathsf T}M_0.
  \]
  Matrix multiplication gives the stated entries and negative value.
  For $e_1=(1,0)^{\mathsf T}$, the first random contribution is
  $\|M_1e_1\|^2=5/4$, whereas $e_1^{\mathsf T}Ke_1=4/3$.
  Thus, it is not valid even as an upper bound for the full coefficient.

  The entire series are absolutely Gaussian integrable, including their
  squares, by the bound in No-go~\ref{P:nogo-14} with the additional
  polynomial factor from $zs(z)$. Absolute summation of the even-degree
  terms at a fixed nonzero $\sigma_1$ bounds the tail of degree at least
  four by $C\sigma^4$ for $|\sigma|\le|\sigma_1|$. Therefore, for
  $x=(1,-2)^{\mathsf T}$ the actual expected energy satisfies
  \[
    \E\|M(\sigma Z)x\|^2=5-\frac{16}{45}\sigma^2+O(\sigma^4).
  \]
  This proves No-go~\ref{P:nogo-15}.
\end{proof}

The example refutes arbitrary-phase covariance monotonicity for this
scalar flow. It says nothing against monotonicity of the original
solution-value moment with initial data $(1,0)$ in the homogeneous model.
Section~\ref{S:static-phase} retains both contributions to the coefficient.
The all-order estimate of Proposition~\ref{P:static-phase-all-orders} and
the interaction bound of Corollary~\ref{C:static-interaction} retain these
terms. Positivity of the first random variance does not control the cross
term with the mean.

\subsection{NG-16: replacing a tilted mean by a deterministic value}
\label{S:nogo-16}

The Gaussian density cost is exact. A loss may instead enter when the
random solution under the tilted law is replaced by a deterministic
surrogate. The following two calculations identify the loss without
asserting an obstruction to Gaussian shifts themselves.

\setcounter{nogo}{15}
\begin{nogo}[A deterministic surrogate can lose the annealed factor]
  \label{P:nogo-16}
  Let $A>0$ and $0<r<2$, and define
  \[
    J_c=\sup_{H\ge0}\{cAH^r-H^2/2\}.
  \]
  Then $J_c=c^{2/(2-r)}J_1$, whereas
  \[
    \sup_{H\ge0}\{2AH^r-H^2\}=2J_1<J_2.
  \]
  Moreover, for $Y=e^{\sigma Z}$, $Z$ standard Gaussian and $\sigma>0$,
  the full tilted-mean Cauchy--Schwarz bound recovers $\E Y^2$ exactly,
  while substituting $Y(h)=e^{\sigma h}$ loses a factor in its logarithm.
\end{nogo}
\begin{proof}
  Substitute $H=c^{1/(2-r)}K$ in the definition of $J_c$. Both terms
  acquire the factor $c^{2/(2-r)}$. The supremum $J_1$ is finite, since
  $r<2$, and positive, since for small positive $H$ the positive term
  dominates. Factoring two from the surrogate score proves the strict
  comparison because $2/(2-r)>1$.

  Let $Q_h$ be the Gaussian shift with density
  $L_h=e^{hZ-h^2/2}$ relative to the standard Gaussian law. Then
  $\E L_h^2=e^{h^2}$, so Cauchy--Schwarz gives
  \[
    \E Y^2\ge(\E_{Q_h}Y)^2e^{-h^2}.
  \]
  The exact tilted mean is $\E_{Q_h}Y=e^{\sigma^2/2+\sigma h}$.
  Maximizing its resulting exponent
  $\sigma^2+2\sigma h-h^2$ at $h=\sigma$ gives
  $e^{2\sigma^2}=\E Y^2$. Substitution of the deterministic value
  $Y(h)$ instead gives the maximum $e^{\sigma^2}$.
  This proves No-go~\ref{P:nogo-16}.
\end{proof}

These are generic scalar score calculations, not an SPDE no-go. They
identify which term is lost and show that retaining the full tilted mean
can avoid the loss. A deterministic surrogate remains useful if one
proves that its error is negligible on the intended exponential scale;
without such an estimate it cannot identify the annealed constant.

\subsection{NG-20: covariance profiles need not separate}
\label{S:nogo-20}

A fixed spatial profile times a scalar temporal factor is easier to
propagate than a general spacetime control. Such a product family is
legitimate as a restriction, but it does not represent every positive
control even in the white spatial model.

\setcounter{nogo}{19}
\begin{nogo}[A positive control need not have proportional covariance profiles]
  \label{P:nogo-20}
  Fix $0<a<1$ and one-dimensional spatial covariance
  $\Gamma=c_\Gamma\delta_0$, $c_\Gamma>0$. There is a nonnegative smooth
  compactly supported spacetime test $\psi$ for which
  \[
    h_\sigma(x)=c_\Gamma\int_0^1|\sigma-u|^{-a}\psi(u,x)\ud u,
    \qquad 0\le\sigma\le1,
  \]
  cannot be represented as $b(\sigma)U(x)$ with one fixed spatial $U$.
\end{nogo}
\begin{proof}
  Choose a nonzero nonnegative smooth $f_1$ supported in $(1/8,1/4)$,
  let $f_2(u)=f_1(1-u)$, and choose nonzero nonnegative smooth bumps
  $q_1,q_2$ supported respectively in $(-2,-1)$ and $(1,2)$. Set
  $\psi(u,x)=f_1(u)q_1(x)+f_2(u)q_2(x)$. Put
  \[
    A=\int_0^1u^{-a}f_1(u)\ud u,\qquad
    B=\int_0^1(1-u)^{-a}f_1(u)\ud u.
  \]
  Since $u<1-u$ on the support of $f_1$, we have $A>B>0$.
  Reflection in time gives
  \[
    h_0=c_\Gamma(Aq_1+Bq_2),\qquad
    h_1=c_\Gamma(Bq_1+Aq_2).
  \]
  Choose $x_i$ with $q_i(x_i)>0$. Disjoint supports imply
  \[
    h_0(x_1)h_1(x_2)-h_1(x_1)h_0(x_2)
    =c_\Gamma^2(A^2-B^2)q_1(x_1)q_2(x_2)>0.
  \]
  A representation by one $U$ would make this determinant zero.
  The white covariance has $\alpha=1$, so $1+a<2$ verifies the model
  parameter condition. This proves No-go~\ref{P:nogo-20}.
\end{proof}

This is an exact representation obstruction. It does not prove that
separable controls fail to attain a variational supremum; a theorem
reducing optimizers to a smaller family would be a different statement.
A useful weaker fact survives: because $|\sigma-u|^{-a}\ge1$ on the
unit time square, every nonnegative test has
$h_\sigma\ge\Gamma*\int_0^1\psi(u,\cdot)\ud u$.
This common lower profile does not by itself prove a phase-propagation
comparison. Reopening an exact product reduction requires extra structure
on the control, while an optimizer reduction requires a variational proof.

\subsection{NG-31: initial phase and the backward cone}\label{S:nogo-31}

Two seemingly small generalizations remove essential hypotheses: replacing
the specified initial phase by an arbitrary one, and checking agreement
of spatial cutoffs only on the terminal observation set. Simple
deterministic waves already distinguish these statements.

\setcounter{nogo}{30}
\begin{nogo}[Growth does not give every-phase lower bounds or local cutoff consistency]
  \label{P:nogo-31}
  The growth rate of a wave phase propagator need not be a lower growth
  rate for every nonzero initial phase. Also, two smooth initial-data
  masks that agree on an observation interval can give different free-wave
  displacements throughout that interval at a later time, even when their
  initial displacements agree everywhere.
\end{nogo}
\begin{proof}
  For the first assertion, let
  $B=\left(\begin{smallmatrix}0&1\\1&0\end{smallmatrix}\right)$.
  Its eigenvectors $(1,1)^{\mathsf T}$ and $(1,-1)^{\mathsf T}$ have
  eigenvalues $1$ and $-1$. Hence, $\|e^{tB}\|=e^t$ for $t\ge0$,
  but $e^{tB}(1,-1)^{\mathsf T}=e^{-t}(1,-1)^{\mathsf T}$.

  For the second assertion take the one-dimensional free wave with zero
  initial displacement and nonnegative smooth velocity $v$ supported in
  $(5,6)$, normalized by $\int v=1$. Choose smooth masks $\chi_s,\chi_l$
  between zero and one, with $\chi_s=1$ on $[-2,2]$ and support in
  $[-4,4]$, and $\chi_l=1$ on $[-14,14]$. Both masks agree on the
  observation interval $[-2,2]$. The masked velocities are nevertheless
  $\chi_sv=0$ and $\chi_lv=v$. D'Alembert's formula at time nine gives
  \[
    u_s(9,x)=0,\qquad
    u_l(9,x)=\frac12\int_{x-9}^{x+9}v(y)\ud y=\frac12,
    \qquad -2\le x\le2.
  \]
  Indeed, every interval $[x-9,x+9]$ in question contains $(5,6)$.
  This proves No-go~\ref{P:nogo-31}.
\end{proof}

The first example does not refute a terminal lower bound for the original
phase $(1,0)$. The second requires agreement of both displacement and
velocity on the full backward cone, not only on the observer. Finite
propagation with that guard survives. Thus, a successor must retain the
specified phase, or prove a suitable nonvanishing unstable component, and
must impose cutoff agreement on the relevant cone for both coordinates.

\section{Topology, localization and limits of variational arguments}
\label{S:nogo-topology}

A useful approximation must be small in the topology in which the next
argument uses it. The obstructions below distinguish a local operator from a
Hilbert--Schmidt operator, a fixed spatial cutoff from a growing one, and
pointwise convergence from convergence of variational suprema. Throughout,
$a=\alpha_0\in(0,1)$ and $\Lambda=(1-\Delta)^{1/2}$, except that the positive
number $L$ denotes the temporal feature cutoff in No-go~\ref{P:nogo-32}.

\subsection{The first driver cannot always be Hilbert--Schmidt}
\label{S:nogo-21}

The one-dimensional operator construction suggests trying the same
Hilbert--Schmidt space in higher dimensions. A compact spatial observation
removes infinite volume, but it does not remove the following high-frequency
obstruction. The relevant compatibility is with the given noise, through its
bilinear first-driver tests.

\setcounter{nogo}{20}
\begin{nogo}[A same-noise Hilbert--Schmidt obstruction]
  \label{P:nogo-21}
  Suppose that $d\in\{2,3\}$, $s\in\R$, and $0<T\le1$. Put
  $E_s=H^s(\R^d)\times H^{s-1}(\R^d)$ and let $S(t)$ be the massless wave
  phase group. For every nonzero covariance satisfying Hypothesis~\ref{H:homog},
  there are a nonnegative compact smooth physical coefficient mask $\zeta$ and
  a compact smooth observation cutoff $\chi$, equal to one on a sufficiently
  large ball, for which the localized Gaussian first-driver form has no
  compatible realization in $L^2(\Omega;\operatorname{HS}(E_s,E_s))$.
  Here compatibility means equality with the given $W$ for each fixed smooth
  phase pair, as specified in the proof.
\end{nogo}

\begin{proof}
  First consider a deterministic coefficient $h(u,x)=k(u)H(x)$, where
  $H\in C_c^\infty(\R^d)$ is nonzero, even and nonnegative, and
  $k\in W^{1,\infty}[0,T]$ satisfies $k\ge k_0>0$. Write
  $B(b)(v,w)=(0,bv)$ and
  \[
    A_h=\int_0^T S(-u)B(h(u))S(u)\ud u,\qquad
    S(u,\xi)=\begin{pmatrix}
      \cos(u|\xi|)&\sin(u|\xi|)/|\xi|\\
      -|\xi|\sin(u|\xi|)&\cos(u|\xi|)
    \end{pmatrix}.
  \]
  The sine quotient has its continuous value at zero. The integral is a strong
  operator integral on $E_s$: $S(u)$ is strongly continuous and uniformly
  bounded on compact time intervals, and multiplication by $H$ is bounded on
  $H^s$. For the latter assertion, the weighted Fourier kernel is bounded in
  absolute value by a constant times
  $\langle\xi-\eta\rangle^{|s|}|\widehat H(\xi-\eta)|$, an integrable
  convolution kernel. This proves boundedness for every fixed real $s$.

  Normalize the Fourier coordinates by
  $\operatorname{diag}(\langle\xi\rangle^s,\langle\xi\rangle^{s-1})$.
  Using the unitary Fourier transform, the lower-left block of $A_h$ then has
  kernel
  \[
    K_s(\xi,\eta)=(2\pi)^{-d/2}
    \langle\xi\rangle^{s-1}\langle\eta\rangle^{-s}
    \widehat H(\xi-\eta)
    \int_0^T k(u)\cos(u|\xi|)\cos(u|\eta|)\ud u.
  \]
  This formula follows from both phase coordinates: before applying the weights,
  the full product, with $r=|\xi|$, $\rho=|\eta|$, is
  \[
    \widehat H(\xi-\eta)
    \begin{pmatrix}
      -\dfrac{\sin(ru)}r\cos(\rho u)&
      -\dfrac{\sin(ru)}r\dfrac{\sin(\rho u)}\rho\\
      \cos(ru)\cos(\rho u)&\cos(ru)\dfrac{\sin(\rho u)}\rho
    \end{pmatrix}.
  \]
  All Fourier and time exchanges are justified on Schwartz pairs by the rapid
  decay of $\widehat H$ and then by density.

  Let $J_0=\int_0^T k>0$ and
  $B_0=\|k\|_\infty+T\|k'\|_\infty$. Choose $\epsilon>0$ so that
  $\epsilon\le1$, $\epsilon T\le1$ and
  $\widehat H(z)\ge\widehat H(0)/2>0$ when $|z|\le\epsilon$.
  If $|\xi-\eta|\le\epsilon$, then
  \[
    \int_0^T k(u)\cos((r-\rho)u)\ud u\ge J_0/2,
    \qquad
    \left|\int_0^T k(u)\cos((r+\rho)u)\ud u\right|
    \le B_0/(r+\rho).
  \]
  The second bound is integration by parts in the absolutely continuous
  function $k$. Product-to-sum therefore bounds the time integral in $K_s$
  from below by $J_0/8$ when
  $r\ge R_0:=\max(1,2\epsilon,4B_0/J_0)$.
  On the same region,
  \[
    \langle\xi\rangle^{s-1}\langle\eta\rangle^{-s}
    =\langle\xi\rangle^{-1}
    (\langle\xi\rangle/\langle\eta\rangle)^s
    \ge 2^{-|s|}/(\sqrt2 r).
  \]
  Consequently, $|K_s(\xi,\eta)|\ge c_s/|\xi|$ there, for $c_s>0$.
  If $P_N$ cuts off Fourier frequencies to $B_N$, its truncated kernel is
  square integrable and Parseval followed by Tonelli gives
  \[
    \|P_NK_sP_N\|_{\operatorname{HS}}^2
    \ge c_s^2|B_\epsilon||\mathbb S^{d-1}|
    \int_{R_0}^{N-\epsilon}r^{d-3}\ud r\longrightarrow\infty.
  \]
  Indeed Parseval in the second kernel variable identifies its square integral
  with the sum of squared images of an orthonormal basis. If $A_h$ were
  Hilbert--Schmidt, every block compression would have bounded
  Hilbert--Schmidt norm. The displayed divergence proves otherwise. Real
  Hilbert--Schmidt operators and their complexifications have the same norm,
  so Fourier complexification has not changed the conclusion.

  Suppose that $\operatorname{supp}H\subset B_A$, and let $\chi=1$ on
  $B_{A+T+1}$. With $D_\chi$ denoting multiplication of both phase coordinates,
  finite propagation gives
  \[
    D_\chi A_h=A_h=A_hD_\chi.
  \]
  For completeness, the local wave energy on a shrinking ball has boundary
  flux $v_t\partial_n v-(v_t^2+|\nabla v|^2)/2\le0$. Zero initial data
  therefore remain zero in the interior cone; time reversal gives the same
  support assertion for $S(-u)$. Apply this first to smooth data and then use
  Sobolev continuity and distributional testing to pass to $E_s$. The
  coefficient acts only in $B_A$, and the resulting output lies in
  $B_{A+T}$. Thus the failure persists after this compact observation.

  We next produce $h$ as an actual covariance mean. Choose $L_0\ge1$ with
  $0<\Gamma(B_{L_0})<\infty$, and an even nonnegative $q\in C_c^\infty$
  equal to at least one on $B_{2L_0+1}$. Set $U=\Gamma*q$. Local finiteness
  of $\Gamma$ permits differentiation on every compact set, so $U$ is smooth,
  even and nonnegative, with $U(0)>0$. Moreover,
  \[
    Q=\int q(x)U(x)\ud x
    \ge |B_{L_0+1}|\Gamma(B_{L_0})>0.
  \]
  Choose an even nonnegative compact smooth $\zeta$, equal to one near zero,
  and put $H=\zeta U$. Choose $f\in C_c^\infty(0,1)$ nonnegative and
  nonzero, and set
  \[
    Z=W(f\otimes q),\qquad
    k(u)=\int_0^1|u-v|^{-a}f(v)\ud v.
  \]
  For $u\in[0,1]$, $k(u)\ge\int f>0$. Moving the derivative to the smooth
  zero extension of $f$ shows $k\in C^1[0,1]$, since $|\cdot|^{-a}$ is
  locally integrable. Also $\E Z^2=Q\int_0^1 f(u)k(u)\ud u\in(0,\infty)$.
  The mask is applied after the covariance convolution: $H=\zeta(\Gamma*q)$.

  It remains to specify the compatible first-driver form with the correct
  Sobolev adjoint. For real Schwartz phases $z,w$ put
  \[
    v_u=[S(u)z]_1,\qquad
    G_w(u)=-\frac{\sin(u|D|)}{|D|}(\chi\Lambda^{2s}w_1)
    +\cos(u|D|)(\chi\Lambda^{2s-2}w_2),
  \]
  and define
  $F_{z,w}(u,x)=\mathbf1_{(0,T)}(u)\zeta(x)v_u(x)G_w(u,x)$.
  This test lies in the raw noise completion. In fact, on a fixed compact
  spatial support $K$, its covariance norm is bounded by a constant times
  $\|u\mapsto\|F(u)\|_\infty\|_{L^2(0,T)}$, by local Radon mass
  $\Gamma(K-K)$ and the time Schur bound for $|u-v|^{-a}$. Smooth interior
  time approximation gives membership. Define
  $\ell_\chi(z,w)=W(F_{z,w})$ on the given noise space.
  Direct transposition, with the Sobolev weights retained, gives
  \[
    G_w(u)=\Lambda^{2s-2}
    \big((D_\chi S(-u))^{*_{E_s}}w\big)_2.
  \]
  In particular, multiplication by $\chi$ has not been treated as self-adjoint
  in $E_s$. Absolute Fubini in the covariance pairing now gives
  \[
    \E[Z\ell_\chi(z,w)]
    =\int_0^T\!\int k(u)H(x)v_u(x)G_w(u,x)\ud x\ud u
    =\langle D_\chi A_hz,w\rangle_{E_s}
    =\langle A_hz,w\rangle_{E_s}.
  \]
  All spatial factors in the raw tests have compact support and the temporal
  singularity is integrable, which supplies the absolute Fubini bound.

  If $X\in L^2(\Omega;\operatorname{HS}(E_s,E_s))$ satisfied
  $\langle Xz,w\rangle_{E_s}=\ell_\chi(z,w)$ almost surely for every fixed
  Schwartz pair, then Cauchy--Schwarz would make
  $Y=\E[ZX]$ a Hilbert--Schmidt Bochner integral. Continuous matrix evaluation
  would give $\langle Yz,w\rangle=\langle A_hz,w\rangle$ for all such pairs.
  Density implies $Y=A_h$, a contradiction. This proves No-go~\ref{P:nogo-21}.
\end{proof}

This result concerns the same-index phase spaces $E_s$ and a compatible
$L^2$ Hilbert--Schmidt realization. It does not concern every cutoff: the zero
cutoff, for example, has no such obstruction. It is not a nonexistence result
for the Stratonovich solution or its moments. A route using loss of regularity
between input and output, operator norms, or a different Schatten class must
supply its own estimates; changing the fixed index $s$ alone cannot reopen
this particular route.

\subsection{A physical mask may not act on the noise quotient}
\label{S:nogo-22}

A smooth physical mask is harmless on ordinary test functions. It need not
preserve the zero classes of a covariance seminorm. The following is the
obstruction already proved in Proposition~\ref{P:masked-quotient-obstruction};
we include its witness to explain why the auxiliary completion is needed.

\setcounter{nogo}{21}
\begin{nogo}[A mask that separates noise-null representatives]
  \label{P:nogo-22}
  For $d=2$, $\alpha=1$, $0<a<1$ and
  $\Gamma=\delta_0(\ud x_1)\ud x_2$, there exist real compact smooth $F$ and
  a nonnegative smooth square mask $\zeta$ for which
  $\E W(F)^2=0$ but $\E W(\zeta F)^2>0$.
  Thus multiplication by this mask does not descend to the original noise
  quotient, and no finite uniform covariance-norm multiplier bound holds.
\end{nogo}
\begin{proof}
  Choose a nonnegative nonconstant compact smooth plateau function $\psi$ and
  put $\chi=\psi^2$. The spatial covariance seminorm is
  \[
    V(g)=\int_\R\left(\int_\R g(x_1,x_2)\ud x_2\right)^2\ud x_1.
  \]
  This also shows positive definiteness of the nonnegative even tempered
  measure $\Gamma$; its scaling is $\Gamma(cA)=c\Gamma(A)$, as required
  for $d=2$, $\alpha=1$. Put
  $g(x)=\chi(x_1)\chi''(x_2)$ and
  $\zeta(x)=\chi(x_1)\chi(x_2)=(\psi(x_1)\psi(x_2))^2$.
  Then $\int\chi''=0$, whereas integration by parts gives
  \[
    \int\zeta(x)g(x)\ud x_2=-\chi(x_1)^2\int(\chi')^2.
  \]
  Consequently, $V(g)=0$ and
  $V(\zeta g)=\int\chi^4\,\big(\int(\chi')^2\big)^2>0$.
  Tensor with a nonzero nonnegative $\tau\in C_c^\infty(0,T)$. Its temporal
  variance is finite and strictly positive, so $F=\tau\otimes g$ has the
  asserted full-noise variances. This proves No-go~\ref{P:nogo-22}.
\end{proof}

This covariance is not two-dimensional spatial white noise. One can still
mask individual tests before taking an auxiliary completion. To reopen a
multiplier argument on the original quotient, one must first prove that the
chosen mask preserves its null space and then prove the required norm bound.

\subsection{Frequency decay does not give whole-space compactness}
\label{S:nogo-24}

The compact-cutoff estimates use spatial confinement as well as frequency
smoothing. Removing the former loses compactness even for a deterministic
Bessel multiplier; see Proposition~\ref{P:cutoff-boundary-tests}.

\setcounter{nogo}{23}
\begin{nogo}[Translations defeat whole-space compactness]
  \label{P:nogo-24}
  For every $d\ge1$ and $\gamma>0$, the operator $\Lambda^{-\gamma}$ on
  $L^2(\R^d)$ is not compact.
\end{nogo}
\begin{proof}
  Take a nonzero Schwartz function $f$ with Gaussian Fourier transform and put
  $g=\Lambda^{-\gamma}f$. Then $g$ is nonzero and Schwartz. For translations
  $T_hg(x)=g(x-h)$, Plancherel gives
  \[
    \langle T_hg,g\rangle
    =\int e^{-i\xi\cdot h}|\widehat g(\xi)|^2\ud\xi\longrightarrow0
    \quad (|h|\longrightarrow\infty).
  \]
  The limit follows, for example, by integration by parts in the Fourier
  transform of the Schwartz function $|\widehat g|^2$.
  Choose a sufficiently large spacing $b$ and $h_n=nb e_1$ so that distinct
  images $T_{h_n}g$ have squared distance at least $\|g\|_2^2$.
  The inputs $T_{h_n}f$ have constant norm and
  $\Lambda^{-\gamma}T_{h_n}f=T_{h_n}g$. Their images have no Cauchy
  subsequence. This proves No-go~\ref{P:nogo-24}.
\end{proof}

Compact spatial cutoffs and the conditional Schatten estimates in the main
text survive. No necessity at a Schatten endpoint is asserted here. A
whole-space successor needs a tightness condition that prevents translation
of mass, or an explicitly different topology.

\subsection{Fixed projections do not control growing spatial masks}
\label{S:nogo-25}

At each fixed radius the Cameron mean map can be compact. Its finite-rank
approximations need not be uniform in that radius. The following witness uses
actual unit controls for one-dimensional spatial white noise. The fixed-radius
map and its compact approximation are those of
Proposition~\ref{P:cameron-full-mean}.

\setcounter{nogo}{24}
\begin{nogo}[Escape from every fixed smooth projection]
  \label{P:nogo-25}
  Let $\Gamma=\delta_0$ in $d=1$, and let $H_I$ be the closed raw-noise
  space generated by tests supported in $I\times\R$, $I=(0,1)$. Let $P_J$
  be any fixed finite-rank orthogonal projection in $H_I$ whose range is
  spanned by smooth compact tests supported in $I\times\R$. For the physical masked mean multiplier $T_R$ on
  $F=(H^{1/2}\times H^{-1/2})\oplus\R^2$, with the last coordinates
  representing constant displacement and velocity, one has, for every
  $2<r<2/a$,
  \[
    \sup_R\|T_R(I-P_J)\|_{H_I\to L^r((0,1);\mathcal K(F))}\ge c_a>0,
  \]
  where $c_a$ can be chosen independently of $J$. Here the smooth masks $m_R$
  equal one on $[-R,R]$. For a smooth control $\Psi$ with covariance mean
  $h_\Psi$, the operator sends $(v,w,c,d)$ to
  $(0,m_Rh_\Psi(s)(v+c),0,0)$.
\end{nogo}
\begin{proof}
  Use the temporal equilibrium density $f_a$ of
  Lemma~\ref{L:temporal-equilibrium}, so that
  $\int_0^1|s-t|^{-a}f_a(t)\ud t=1$ and
  $\|f_a\|_{H_a}^2=M_a>0$.
  Choose nonnegative $q\in C_c^\infty(\R)$ with $\|q\|_2=1$, and put
  $q_n(x)=q(x-x_n)$, where $x_n\to\infty$. The controls
  \[
    \Psi_n=\frac{f_a\otimes q_n}{\sqrt{M_a}}
  \]
  have unit raw norm and time-independent covariance mean
  $q_n/\sqrt{M_a}$. Membership follows by $L^2$ approximation of $f_a$ and
  the Schur bound for the temporal kernel.
  The finitely many tests spanning $P_J$ have common compact spatial support.
  For all sufficiently large $n$, spatial whiteness gives
  $P_J\Psi_n=0$ exactly. Choose $R_n$ to contain $\operatorname{supp}q_n$
  in the plateau of $m_{R_n}$.

  The constant-displacement input $e=(0,0,1,0)$ has norm one. The phase dual
  \[
    w_n=(0,\Lambda q_n/\|q\|_{H^{1/2}},0,0)
  \]
  also has norm one. Indeed its velocity norm is $H^{-1/2}$, and translation
  preserves $\|q\|_{H^{1/2}}$. The mean multiplier sends $e$ to the velocity
  $q_n/\sqrt{M_a}$, so the $H^{-1/2}$ inner product gives, for every $s$,
  \[
    \langle T_{R_n}\Psi_n(s)e,w_n\rangle_F
    =\frac{\|q_n\|_2^2}{\sqrt{M_a}\|q\|_{H^{1/2}}}
    =:c_a>0.
  \]
  Integrating over the unit time interval proves the bound for all $r$ stated.
  This proves No-go~\ref{P:nogo-25}.
\end{proof}

\begin{figure}
  \centering
  \begin{tikzpicture}[x=1cm,y=1cm]
    \draw[ham/arrow] (-.5,0)--(10.5,0) node[right] {$x$};
    \fill[hamblue!12] (0,.1) rectangle (2.4,1.1);
    \node[ham/note] at (1.2,1.95)
    {Fixed projection $P_J$\\common spatial support};
    \draw[hamblue,thick] (0,.1) rectangle (2.4,1.1);
    \draw[hamteal,thick] (6.5,.1) .. controls (7.1,.1) and (7.1,1.1) .. (7.5,1.1)
    .. controls (7.9,1.1) and (7.9,.1) .. (8.5,.1);
    \node[ham/note] at (7.5,1.65) {Translated bump $q_n$};
    \draw[ham/arrow] (4,.6)--(5.9,.6);
    \node[ham/note] at (4.95,1.15) {$x_n\to\infty$};
    \draw[hamgold,thick] (-.3,-.3)--(-.3,-.6)--(9.5,-.6)--(9.5,-.3);
    \node[ham/note] at (4.6,-1.1) {Growing plateau $m_{R_n}=1$ retains $q_n$};
    \node[ham/note] at (4.6,-1.8)
    {$P_J\Psi_n=0$,\qquad $\|T_{R_n}\Psi_n\|\ge c_a>0$};
  \end{tikzpicture}
  \caption[Escape from a fixed spatial projection]{Schematic for
    No-go~\ref{P:nogo-25}. For one-dimensional spatial white noise, a translated
    unit control is exactly orthogonal to the finitely many compact tests of
    $P_J$, while a growing mask retains its nonzero phase functional.
    At any fixed radius sufficiently distant bumps miss the mask.}
  \label{Fig:nogo-projection-escape}
\end{figure}


For a fixed $R$, sufficiently distant translates miss its mask, so there is
no contradiction with fixed-radius compactness. A replacement must use
projections depending on the radius with quantitative control, or a different
localization argument. The witness does not exclude every such strategy.

\subsection{Two distinct limitations of block costs}
\label{S:nogo-26}

A nonlinear block argument must handle both small perturbations and rare large
blocks. A power larger than one misses the first effect; an averaged budget
alone does not bound the largest block.

\setcounter{nogo}{25}
\begin{nogo}[A pure Young power and a maximum from an average]
  \label{P:nogo-26}
  On $[0,1]$, with $\theta=2/3$ and $\nu=1/\theta=3/2$, no finite constant
  $C>0$ gives
  $\log U(1)\le C\|D\|_{C^\theta}^{\nu}$ for all scalar paths starting at zero
  and solving $\ud U=U\ud D$, $U(0)=1$, when
  $\|D\|_{C^\theta}=\|D\|_\infty+[D]_{C^\theta}$.
  Moreover, a bound on $N^{-1}\sum_{j=1}^N b_j^r$, $r>0$, does not give a
  bound independent of $N$ on $\max_j b_j$.
\end{nogo}
\begin{proof}
  For $D(t)=xt$, $x>0$, the solution is $U(t)=e^{xt}$ and
  $\|D\|_{C^{2/3}}=2x$. Taking $x=1/(16C^2)$ gives
  $C(2x)^{3/2}=x/\sqrt2<x=\log U(1)$.
  For the second assertion, take $b_1=N^{1/r}$ and $b_j=0$ for $j>1$.
  Their averaged $r$th power is one while their maximum tends to infinity.
  This proves No-go~\ref{P:nogo-26}.
\end{proof}

A linear-plus-power cost is not refuted. For any $\tau>0$ and $\nu>1$,
$x\le\tau+\tau^{1-\nu}x^\nu$ gives an explicit alternative with slack.
Also, if $0<\nu<r$ and $N^{-1}\sum b_j^r\le M$, then H\"older gives
$N^{-1}\sum b_j^\nu\le M^{\nu/r}$, and
\[
  \frac1N\sum_{b_j>K}b_j^\nu\le K^{\nu-r}M.
\]
Thus an averaged tail argument survives. Reopening the block route requires
an actual evolution estimate with its linear term or slack, followed by an
averaged bad-block estimate. These scalar tests alone do not construct a
Cameron potential or prove a wave evolution bound.

\subsection{Concentrating temporal controls restrict integrability}
\label{S:nogo-27}

Compactness of a temporal mean map has a natural upper endpoint. A
concentrating control detects both excessive integrability and failure of
compactness at that endpoint.

\setcounter{nogo}{26}
\begin{nogo}[The temporal concentration threshold]
  \label{P:nogo-27}
  Let $I=(0,1)$, $0<a<1$, and let $H_a$ be the completion of smooth tests for
  $\|f\|_{H_a}^2=\int_I\int_I f(t)f(s)|t-s|^{-a}\ud t\ud s$.
  On smooth tests set $C_af(t)=\int_I|t-s|^{-a}f(s)\ud s$.
  The map $C_a$ is not bounded from $H_a$ to $L^r(I)$ when $r>2/a$,
  including $r=\infty$, and it is not compact into $L^{2/a}(I)$.
\end{nogo}
\begin{proof}
  Choose $t_0\in I$ and a nonzero nonnegative
  $\phi\in C_c^\infty(-1,1)$. Put
  \[
    v_\phi=\int\!\int\phi(u)\phi(v)|u-v|^{-a}\ud u\ud v\in(0,\infty),
    \qquad
    f_\epsilon(t)=\frac{\epsilon^{a/2-1}}{\sqrt{v_\phi}}
    \phi((t-t_0)/\epsilon).
  \]
  For sufficiently small $\epsilon$, its support lies in $I$. Changing both
  variables shows $\|f_\epsilon\|_{H_a}=1$, and
  \[
    (C_af_\epsilon)(t_0+\epsilon y)
    =\frac{\epsilon^{-a/2}}{\sqrt{v_\phi}}
    \int|y-v|^{-a}\phi(v)\ud v.
  \]
  For $|y|\le1$, the last integral is at least $2^{-a}\int\phi>0$.
  Thus for finite $r$,
  $\|C_af_\epsilon\|_{L^r}\ge c_r\epsilon^{1/r-a/2}$, and the analogous
  supremum lower bound is $c\epsilon^{-a/2}$. This proves unboundedness above
  $2/a$.
  For every fixed $t\ne t_0$, on the other hand,
  $C_af_\epsilon(t)=O_t(\epsilon^{a/2})\to0$.
  At $r=2/a$ the norm lower bound stays strictly positive. If a subsequence
  converged strongly in this $L^r$ space, a further subsequence would converge
  almost everywhere to its strong limit. The limit would have to be zero,
  contradicting the norm lower bound. This proves No-go~\ref{P:nogo-27}.
\end{proof}

The endpoint argument refutes compactness, not boundedness. Tensoring the
controls with one fixed nonzero spatial profile transfers this input
obstruction to a corresponding separated mean map. The corrected compact
range $2<r<2/a$ is consistent with the witness. A larger range requires
additional temporal regularity or restrictions that exclude concentration.

\subsection{Pointwise convergence does not pass through a supremum}
\label{S:nogo-30}

The Gaussian action cost must remain in a variational comparison. Even with
that cost, bounded near-maximizers can escape in infinitely many directions.

\setcounter{nogo}{29}
\begin{nogo}[Escaping variational scores]
  \label{P:nogo-30}
  There are continuous functions $J_n$ on the Hilbert space $\ell^2$ tending
  to zero at every point such that
  \[
    \sup_h\{J_n(h)-\tfrac12\|h\|^2\}\ge\tfrac12,
    \qquad
    \sup_h\{0-\tfrac12\|h\|^2\}=0.
  \]
  The witnesses for the first inequality have norm one.
\end{nogo}
\begin{proof}
  Let $(e_n)$ be the standard orthonormal sequence and set
  $J_n(h)=(1-2\|h-e_n\|)_+$.
  For each fixed $h\in\ell^2$, its $n$th coordinate tends to zero, and hence
  $\|h-e_n\|^2=\|h\|^2+1-2\langle h,e_n\rangle\to\|h\|^2+1\ge1$.
  Thus $J_n(h)=0$ eventually. But $J_n(e_n)=1$, so its score at $e_n$ is
  $1/2$. The limiting score is nonpositive and equals zero at $h=0$.
  This proves No-go~\ref{P:nogo-30}.
\end{proof}

A bounded action budget alone supplies no strong compactness in $\ell^2$.
To pass to an exact constant one needs a variational upper-limit argument,
for example compactness of near-maximizers in a topology with the needed
upper semicontinuity. Pointwise convergence, or compactness proved only with
an additional fixed spatial cutoff, does not supply that argument by itself.

\subsection{A cutoff point feature can lie outside the raw-noise image}
\label{S:nogo-32}

Lemma~\ref{L:source-retaining-isometry} embeds the original noise into one
fixed larger Gaussian space. Its closed image retains all original raw tests.
It does not contain every positive-cutoff point feature. The following
explicit orthogonal vector explains the role of the auxiliary noise.

\setcounter{nogo}{31}
\begin{nogo}[Positive-cutoff point sources need auxiliary noise]
  \label{P:nogo-32}
  Suppose that $\Gamma\ne0$ satisfies Hypothesis~\ref{H:homog}, $t_0>0$,
  $L\in(0,\infty)$ and $\epsilon>0$. In the feature space
  \[
    \mathcal G=L^2((0,\infty)_\lambda\times\R_r,
    \nu_a(\ud\lambda)\ud r;H_\Gamma),\qquad
    \nu_a(\ud\lambda)=\lambda^{a-1}\ud\lambda/\Gamma_E(a),
  \]
  let $K=\Phi H_W$ be the closed image of
  Lemma~\ref{L:source-retaining-isometry}. Set $g=p_\epsilon(x-\cdot)$ and
  \[
    v(\lambda,r)=\mathbf1_{\{\lambda\le L\}}\sqrt{2\lambda}
    e^{-\lambda(t_0-r)}\mathbf1_{\{r\le t_0\}}g.
  \]
  Then $v\notin K$. On the fixed extension
  $\mathcal I(v)=W(\Phi^{-1}Pv)+J((I-P)v)$, where $P$ projects onto $K$
  and $J$ is independent of $W$, one has
  $\operatorname{Var}(\mathcal I(v)\mid\sigma(W))>0$.
\end{nogo}
\begin{proof}
  The heat profile has finite strictly positive norm:
  \[
    \|g\|_{H_\Gamma}^2=\int p_{2\epsilon}(y)\Gamma(\ud y)>0.
  \]
  Finiteness follows from temperedness and Gaussian decay, and positivity from
  nonzero nonnegative $\Gamma$. Integrating in $r$ also gives
  $\|v\|_{\mathcal G}^2=\|g\|^2\nu_a((0,L])<\infty$.
  Choose nonnegative nonzero smooth bumps $u_A,u_B$ supported in compact
  intervals $A\subset(0,L)$ and $B\subset(L,\infty)$, respectively, and put
  \[
    p(\lambda)=\frac{u_A(\lambda)}{\int u_A\ud\nu_a}
    -\frac{u_B(\lambda)}{\int u_B\ud\nu_a}.
  \]
  Then $\int p\ud\nu_a=0$ and $\int_{(0,L]}p\ud\nu_a=1$.
  Choose $\varphi\in C_c^\infty(0,\infty)$ with $\varphi(t_0)=1$, and define
  \[
    q(\lambda,r)=\frac{p(\lambda)}{\sqrt{2\lambda}}
    (\varphi'(r)+\lambda\varphi(r))g.
  \]
  Its support in $\lambda$ stays away from zero and infinity. In particular,
  \[
    \|q\|_{\mathcal G}^2=\frac{\|g\|^2}{2}
    \left(\|\varphi'\|_2^2\int\frac{p^2}{\lambda}\ud\nu_a
    +\|\varphi\|_2^2\int\lambda p^2\ud\nu_a\right)<\infty,
  \]
  where the cross term vanishes because $\int\varphi'\varphi=0$.
  For a smooth compact raw tensor $F$ the isometry is
  \[
    \Phi F(\lambda,r)=\sqrt{2\lambda}
    \int_{t>0,\ t\ge r}e^{-\lambda(t-r)}F_t\ud t.
  \]
  The identity
  \[
    \int_{-\infty}^t e^{-\lambda(t-r)}
    (\varphi'(r)+\lambda\varphi(r))\ud r=\varphi(t)
  \]
  follows by differentiating $e^{-\lambda(t-r)}\varphi(r)$ in $r$.
  All integrals paired with $q$ are absolutely integrable because its two
  supports are compact and its $\lambda$ support stays away from zero.
  Consequently,
  \[
    \langle\Phi F,q\rangle_{\mathcal G}
    =\left(\int_{t>0}\langle F_t,g\rangle_{H_\Gamma}
    \varphi(t)\ud t\right)\int p\ud\nu_a=0.
  \]
  Density and continuity give $q\perp K$. The same calculation for the point
  feature instead gives
  \[
    \langle v,q\rangle_{\mathcal G}
    =\varphi(t_0)\|g\|^2\int_{(0,L]}p\ud\nu_a=\|g\|^2>0.
  \]
  Thus $(I-P)v\ne0$. Independence of $J$ and $W$ yields
  \[
    \operatorname{Var}(\mathcal I(v)\mid\sigma(W))
    =\|(I-P)v\|^2\ge\frac{|\langle v,q\rangle|^2}{\|q\|^2}>0.
  \]
  This proves No-go~\ref{P:nogo-32}.
\end{proof}

\begin{figure}
  \centering
  \begin{tikzpicture}[x=1cm,y=1cm]
    \filldraw[fill=hamblue!7,draw=hamblue] (-.5,-.5)--(6,-.5)--(7.5,1.4)--(1,1.4)--cycle;
    \coordinate (o) at (1,.15);
    \coordinate (p) at (4.5,.65);
    \coordinate (v) at (4.5,3.05);
    \draw[ham/arrow] (o)--(p) node[midway,below] {$Pv$};
    \draw[ham/arrow] (o)--(v) node[midway,above left] {$v$};
    \draw[ham/arrow,draw=hamteal] (p)--(v)
    node[ham/note,midway,right,align=left]
    {$(I-P)v\ne0$\\supplied by $J$};
    \draw[hamblue] (4.22,.61)--(4.22,.88)--(4.5,.92);
    \node[ham/note] at (2.5,-.95) {Closed raw-noise image $K=\Phi H_W$};
    \node[ham/note] at (1.45,2.55) {Feature space $\mathcal G$};
  \end{tikzpicture}
  \caption[The point feature and the closed noise image]{Hilbert-space
    schematic for No-go~\ref{P:nogo-32}, drawn in three coordinates only.
    The explicit vector $q\perp K$ has $\langle v,q\rangle>0$, so a
    positive finite cutoff point feature at $t_0>0$ has a nonzero component
    supplied by $J$. This says nothing by itself about measurability of a
    nonlinear limit with respect to $W$.}
  \label{Fig:nogo-feature-range}
\end{figure}


The conclusion concerns positive finite cutoffs and $t_0>0$ in this fixed
extension. It neither treats time zero nor proves that a nonlinear canonical
limit retains an auxiliary component. Raw-test convergence can remove such
components in a limit. The fixed-extension source construction remains valid;
a claim that a particular nonlinear limit is measurable with respect to $W$
alone would need its own convergence argument.

\section{Finite optimization and checks on review arguments}
\label{S:nogo-optimization}

The last group concerns finite calculations. These examples are small enough
that every number can be checked directly, but the distinctions they expose
matter in the infinite-dimensional proof: the two variance directions, the
normalization of an energy, the active set of an optimization, and the exact
inequality used in a review.

\subsection{NG-23: both Gaussian variance directions}
\label{S:nogo-23}

A rectangular Gaussian matrix has an input variance and an output variance.
Keeping only one would simplify the moment estimate, but the dimensions of
the two spaces cannot be ignored. This is the counterexample behind
Remark~\ref{R:gaussian-one-variance-test}.

\setcounter{nogo}{22}
\begin{nogo}[A one-sided variance bound fails]\label{P:nogo-23}
  Let $g_1,\ldots,g_n$ be independent standard real Gaussians,
  $A_i=e_i^{\mathsf T}$ the coordinate row matrices, and
  $X=\sum_i g_iA_i$. Put $V_{\rm in}=\sum_iA_i^*A_i$ and
  $V_{\rm out}=\sum_iA_iA_i^*$. For $n=16$,
  \[
    \E\|X\|_{S_4}^4=288>144=9\operatorname{Tr}V_{\rm in}^2.
  \]
  The transposed family gives the analogous failure of an output-only bound.
  More generally, no constant independent of $n$ bounds the left side by
  $\operatorname{Tr}V_{\rm in}^2$ alone for all these row matrices.
\end{nogo}

\begin{proof}
  The only nonzero singular value of $X$ is
  $(\sum_i g_i^2)^{1/2}$. Since $\E g_i^4=3$ and
  $\E g_i^2g_j^2=1$ for $i\ne j$,
  \[
    \E\|X\|_{S_4}^4
    =\E\left(\sum_i g_i^2\right)^2
    =3n+n(n-1)=n(n+2).
  \]
  Here $V_{\rm in}=I_n$, whereas $V_{\rm out}=n$ acts on a
  one-dimensional space. The proposed bound gives $9n$, which is smaller
  than $n(n+2)$ at $n=16$. The ratio to $\operatorname{Tr}V_{\rm in}^2=n$
  is $n+2$, so no fixed constant works for all $n$. Transposition exchanges
  the variance directions and preserves the singular values. This proves
  No-go~\ref{P:nogo-23}.
\end{proof}

\paragraph{Scope and what survives.}
This is a finite-matrix counterexample, not a counterexample to a wave
estimate. Proposition~\ref{P:finite-gaussian-matrix} retains both variances.
At $n=16$, its average and maximum bounds have fourth powers $1224$ and
$2304$, respectively, and neither is contradicted by $288$.

\paragraph{What would permit a different estimate?}
A one-sided estimate would need additional structure relating the two
variance operators, with its dimension dependence stated. It cannot follow
from the Gaussian matrix form alone.

\subsection{NG-28: the normalization matters when bounds are compared}
\label{S:nogo-28}

The occupation argument gives the upper ceiling $U_2$ for $C_2$.
It is tempting to compare it with an energy expression by changing the
coefficient of the gradient penalty. The comparison must be calculated:
a value with penalty one is not interchangeable with a value with penalty
one-half. In particular, not every comparison value below is itself a proved
upper bound for the moment constant.

Retain $q=4-\alpha$, $k=3-\alpha$ and $m=q/k$, and the occupation energies
$\mathcal{E}_\zeta^{\mathcal O}$ of~\eqref{E:raw-occupation-definitions}.
For $z>0$, the scalar maps are
\[
  F_b(z)=k\left[
  \frac{(q-2b)^{(q-2b)/2}}{8(q-b)^{q-b}}
  \left(\frac{2z}{2-\alpha}\right)^{(2-\alpha)/2}
  \right]^{1/k},\qquad 0\le b<q/2.
\]
These are the maps in~\eqref{E:occupation-scalar-map}; here $b$ is only
a display parameter and $a=\alpha_0$ is the temporal exponent.

\setcounter{nogo}{27}
\begin{nogo}[An occupation ceiling need not improve a comparison value]
  \label{P:nogo-28}
  For $d=2$, $\alpha=5/4$, $a=1/2$ and every nonzero admissible covariance,
  \[
    2^mF_a(\mathcal{E}_1^{\mathcal O})<U_2.
  \]
  Thus, the reverse comparison is not valid throughout the admissible class.
\end{nogo}

\begin{proof}
  The energies are finite and positive by
  Proposition~\ref{P:raw-occupation-energy}. Put
  \[
    \Delta=\frac12\{q\log q+(q-2a)\log(q-2a)\}
    -(q-a)\log(q-a).
  \]
  Subtracting the logarithms of $F_a$ and $F_0$ gives
  $\log(F_a(z)/F_0(z))=\Delta/k$. The identity in
  Theorem~\ref{T:intrinsic-occupation-ceiling} gives
  $U_2=2^m(2\zeta)^{\alpha/(2k)}F_0(\mathcal{E}_\zeta^{\mathcal O})$.
  Therefore,
  \[
    \frac{2^mF_a(\mathcal{E}_\zeta^{\mathcal O})}{U_2}
    =\exp\left\{\frac{\Delta-(\alpha/2)\log(2\zeta)}{k}\right\}.
  \]
  At the stated parameters, $q=11/4$, $k=7/4$, and
  \[
    8\{\Delta-(\alpha/2)\log2\}
    =\log\frac{11^{11}7^7}{2^5 9^{18}}<0,
  \]
  because
  \[
    234966429149994773<4803028329503971872.
  \]
  This proves No-go~\ref{P:nogo-28}.
\end{proof}

\paragraph{Scope and what survives.}
This refutes an ordering of two computed quantities, not an inequality for
the actual moment. It does not decide whether $C_2$ is below the smaller
comparison value. At $\zeta=1/2$, the logarithmic difference is instead
$\Delta/k>0$, by strict convexity of $x\log x$; that is the valid strict
comparison in Corollary~\ref{C:intrinsic-spatial-bracket}.

\paragraph{What would permit a different comparison?}
Match the covariance amplitude, admissible profiles and kinetic penalty,
then prove the sign of the resulting ratio in the parameter range claimed.
Comparing the names of the energies is insufficient.

\subsection{NG-29: three safeguards in the Gram optimization}
\label{S:nogo-29}

The finite optimization of Section~\ref{S:finite-positive-gram} uses
\[
  F(\Phi)=2\sum_i\ell_i
  \sqrt{(\langle\Phi,k_i\rangle-b_i)_+}-\frac12\|\Phi\|^2,
  \qquad \ell_i>0,\quad b_i\ge0.
\]
The Gram entries $C_{ij}=\langle k_i,k_j\rangle$ are nonnegative. Three
possible simplifications fail: omitting the empty active set, invoking
global concavity, and treating entrywise order as positive semidefinite order.

\setcounter{nogo}{28}
\begin{nogo}[Empty sets, concavity and matrix order]\label{P:nogo-29}
  Each of the following failures occurs within the finite setting above.
  \begin{enumerate}
    \item With one unit vector, $\ell=1/10$ and $b=4$, the nonempty-set
      dual infimum is negative, while the primal supremum is nonnegative.
    \item With $H=\R$ and $k=\ell=b=1$, the primal objective is not concave.
    \item There are positive semidefinite Grams $C,D$ with
      $0\le C_{ij}\le D_{ij}$ for every entry but $D-C$ not positive semidefinite.
  \end{enumerate}
\end{nogo}

\begin{proof}
  In the first example the nonempty dual is
  \[
    D(w)=\frac{1}{100w}-4w+\frac{w^2}{2},\qquad w>0.
  \]
  Its infimum is at most $D(4)=-3199/400<0$, whereas $F(0)=0$.
  Dropping the empty value from the maximum of dual infima therefore cannot
  give the primal supremum.

  In the second example $F(x)=2\sqrt{(x-1)_+}-x^2/2$. Directly,
  $F(0)=F(2)=0$ but $F(1)=-1/2$, violating midpoint concavity.

  For the third example take
  \[
    C=\begin{pmatrix}1&0\\0&1\end{pmatrix},\qquad
    D=\begin{pmatrix}1&1/2\\1/2&1\end{pmatrix}.
  \]
  The eigenvalues of $D$ are $3/2$ and $1/2$, while those of $D-C$ are
  $1/2$ and $-1/2$. All entries are ordered as claimed. This proves
  No-go~\ref{P:nogo-29}.
\end{proof}

\paragraph{Scope and what survives.}
These are the checks in Remark~\ref{R:gram-controls}. They explain the
empty-set convention and active-set proof of
Proposition~\ref{P:gram-positive-optimizer}. Entrywise comparison still
works in Proposition~\ref{P:gram-entrywise-comparison}: for nonnegative
weights $w_i$, the form $\sum_{i,j}w_i(D_{ij}-C_{ij})w_j$ is nonnegative.
No order on all vectors, and no covariance monotonicity for nonlinear waves,
is asserted by that calculation.

\paragraph{What would permit a shortcut?}
A restricted concavity argument must first justify its active region.
A matrix comparison must specify the vectors on which its quadratic form
is nonnegative. Neither step follows from a convenient change of terminology.

\subsection{NG-33: check the mathematics returned by a reviewer}
\label{S:nogo-33}

A review report is also a mathematical argument and can contain an incorrect
step. This historical record groups three such checks. The manuscript had
the correct expressions; the examples below explain why the proposed review
steps could not replace them. These are not three new counterexamples to the
stochastic model.

\setcounter{nogo}{32}
\begin{nogo}[Three invalid proof steps]\label{P:nogo-33}
  The following assertions fail.
  \begin{enumerate}
    \item The free phase matrix
      $S_0(t)=\left(\begin{smallmatrix}1&t\\0&1\end{smallmatrix}\right)$
      is a contraction in the Euclidean norm for every $t\ge0$.
    \item For $0<\vartheta<1$, the series term $x^n/(n!)^\vartheta$ equals
      $A_nB_n$, where
      \[
        A_n=\frac{(2x^{1/\vartheta})^n}{n!},\qquad
        B_n=2^{-\vartheta n/(1-\vartheta)}.
      \]
    \item The denominator $k!\,k^{(2-\alpha)k}$ may always be replaced by
      the larger denominator $(k!)^3$ in the coefficient bound of
      Proposition~\ref{P:low-alpha-upper}.
  \end{enumerate}
\end{nogo}

\begin{proof}
  The vector $(0,1)$ has norm one, but its image under $S_0(t)$ has norm
  $\sqrt{1+t^2}>1$ for $t>0$. To transfer this to the whole-space phase,
  use exactly the inhomogeneous Sobolev norm of
  Section~\ref{S:wave-variance}: $E=H^{1/2}\times H^{-1/2}$ and
  $J_0=\operatorname{diag}(\Lambda^{1/2},\Lambda^{-1/2})$, where
  $\Lambda=(1-\Delta)^{1/2}$. With $r=|\xi|$ and
  $\langle\xi\rangle=\sqrt{1+r^2}$, the Fourier symbol of
  $J_0S(t)J_0^{-1}$ is
  \[
    \begin{pmatrix}
      \cos(tr)&\langle\xi\rangle\sin(tr)/r\\
      -r\sin(tr)/\langle\xi\rangle&\cos(tr)
    \end{pmatrix}
    \longrightarrow S_0(t)\qquad(r\downarrow0).
  \]
  Choose normalized Fourier bumps supported in balls shrinking to zero
  and place them in the second coordinate of $J_0z$. Continuity of this
  matrix at zero shows that their image norms tend to $\sqrt{1+t^2}$.
  Thus $S(t)$ is not a contraction on this $E$. The distinction between
  $\langle\xi\rangle$ and $|\xi|$ matters: this is not a claim about the
  homogeneous energy norm. The valid estimate keeps a free-growth factor.

  The correct identity in the second assertion is
  $x^n/(n!)^\vartheta=A_n^\vartheta B_n^{1-\vartheta}$.
  At $\vartheta=3/4$, $n=1$ and $x=1$, its left side is one, whereas
  $A_1B_1=2\cdot(1/8)=1/4$. H\"older applies to the weighted product, as
  in the proof of Proposition~\ref{P:static-phase-all-orders}.

  For the third assertion choose the permitted spatial-white value
  $\alpha=1$ and $k=3$. Then $k!k^{(2-\alpha)k}=162<216=(k!)^3$.
  An upper bound $M/162$ does not imply the smaller upper bound $M/216$.
  The valid factorial weakening is instead
  \[
    k!k^{(2-\alpha)k}\ge(k!)^{3-\alpha},
  \]
  since $k!\le k^k$ and $2-\alpha>0$; at this example its right side is
  $36$. This proves No-go~\ref{P:nogo-33}.
\end{proof}

\paragraph{Scope and what survives.}
The last arithmetic check refutes an inference between numerical bounds;
it does not prove that the actual coefficient exceeds $M/216$.
All three corrected arguments remain in the main text. Their lesson is to
check an asserted derivation against its domains, factors and quantifiers,
even when the report containing it has a favorable overall verdict.

\paragraph{What would permit the discarded steps?}
Contractivity needs a suitable invariant norm or a restricted state class.
The H\"older factors need their conjugate powers. A stronger factorial
bound needs an additional estimate; it cannot be obtained by changing the
exponent in the denominator.

